%%%%%%%%%%%%%%%%%%%%%%%%%%%%%%%%%%%%%%%%%%%%%%%%%%%%%%%%%%%%
% 13.5 FLUID DYNAMICS
% The Composition of Advanced Mathematics
%%%%%%%%%%%%%%%%%%%%%%%%%%%%%%%%%%%%%%%%%%%%%%%%%%%%%%%%%%%%

\section{Fluid Dynamics}
\label{sec:fluid-dynamics}

\prerequisite{
Vector calculus, differential equations, partial differential equations,
continuum mechanics, tensor notation, mathematical physics,
and mathematical modeling.
}

\learningobjectives{
By the end of this section, the reader should be able to:
\begin{itemize}
    \item explain the continuum description of a fluid;
    \item distinguish Eulerian and Lagrangian fluid descriptions;
    \item define velocity, pressure, density, and viscosity fields;
    \item understand streamlines, pathlines, and streaklines;
    \item compute the material derivative;
    \item understand fluid acceleration;
    \item formulate conservation of mass;
    \item derive the continuity equation;
    \item understand incompressible flow;
    \item distinguish compressible and incompressible fluids;
    \item understand pressure in a fluid;
    \item derive hydrostatic pressure;
    \item understand the Cauchy stress tensor for fluids;
    \item distinguish inviscid and viscous fluids;
    \item understand Newtonian fluids;
    \item formulate the Newtonian constitutive stress law;
    \item derive the Euler equations;
    \item derive the Navier--Stokes equations conceptually;
    \item recognize the incompressible Navier--Stokes equations;
    \item understand body forces;
    \item interpret the nonlinear convective term;
    \item recognize steady and unsteady flow;
    \item understand rotational and irrotational flow;
    \item define vorticity;
    \item understand circulation;
    \item recognize Kelvin's circulation theorem conceptually;
    \item understand stream functions in two-dimensional incompressible flow;
    \item understand velocity potentials;
    \item formulate potential flow;
    \item derive Bernoulli's equation under appropriate assumptions;
    \item apply Bernoulli's equation;
    \item understand stagnation pressure;
    \item recognize pressure-velocity conversion;
    \item understand boundary conditions in fluid dynamics;
    \item recognize no-slip and no-penetration conditions;
    \item understand free-surface conditions conceptually;
    \item distinguish internal and external flows;
    \item understand pipe flow;
    \item derive Poiseuille flow;
    \item understand volume flow rate;
    \item recognize shear stress in viscous flow;
    \item understand Reynolds number;
    \item distinguish laminar and turbulent regimes conceptually;
    \item understand nondimensionalization of Navier--Stokes;
    \item recognize inertial and viscous dominance;
    \item understand boundary layers conceptually;
    \item recognize boundary-layer thickness;
    \item understand flow separation conceptually;
    \item recognize drag and lift;
    \item understand dimensional drag coefficients;
    \item recognize potential-flow limitations;
    \item understand circulation-generated lift conceptually;
    \item recognize compressible flow;
    \item define Mach number;
    \item understand speed of sound conceptually;
    \item recognize subsonic, transonic, supersonic, and hypersonic regimes;
    \item understand shock waves conceptually;
    \item recognize conservation laws across discontinuities;
    \item understand shallow-water models conceptually;
    \item recognize Froude number;
    \item understand gravity waves;
    \item recognize geophysical fluid dynamics conceptually;
    \item understand Coriolis effects conceptually;
    \item recognize rotating-fluid balances;
    \item understand diffusion and advection of scalars;
    \item formulate advection-diffusion equations;
    \item recognize Peclet number;
    \item understand thermal convection conceptually;
    \item recognize buoyancy-driven flow;
    \item understand Rayleigh number conceptually;
    \item recognize turbulence modeling challenges;
    \item understand Reynolds averaging conceptually;
    \item recognize Reynolds stresses;
    \item understand numerical fluid simulation conceptually;
    \item distinguish finite difference, finite volume, and finite element
          viewpoints conceptually;
    \item recognize conservation requirements in numerical schemes;
    \item understand CFL-type time-step restrictions conceptually;
    \item recognize model and numerical uncertainty;
    \item connect fluid dynamics with continuum mechanics, PDEs,
          numerical mathematics, transport, control, and physics.
\end{itemize}
}

%%%%%%%%%%%%%%%%%%%%%%%%%%%%%%%%%%%%%%%%%%%%%%%%%%%%%%%%%%%%
\subsection{What Is Fluid Dynamics?}
\label{subsec:what-is-fluid-dynamics}
%%%%%%%%%%%%%%%%%%%%%%%%%%%%%%%%%%%%%%%%%%%%%%%%%%%%%%%%%%%%

Fluid dynamics studies the motion of liquids and gases.

A fluid is modeled as a continuum whose state is described by fields such as

\[
\boxed{
\rho(\mathbf{x},t)
}
\]

for density,

\[
\boxed{
\mathbf{v}(\mathbf{x},t)
}
\]

for velocity,

\[
\boxed{
p(\mathbf{x},t)
}
\]

for pressure, and possibly

\[
\boxed{
T(\mathbf{x},t)
}
\]

for temperature.

The governing equations arise from conservation laws together with
constitutive relations.

%%%%%%%%%%%%%%%%%%%%%%%%%%%%%%%%%%%%%%%%%%%%%%%%%%%%%%%%%%%%
\subsection{Fluids as Continua}
\label{subsec:fluids-as-continua}
%%%%%%%%%%%%%%%%%%%%%%%%%%%%%%%%%%%%%%%%%%%%%%%%%%%%%%%%%%%%

At molecular scales, fluids consist of discrete particles.

At larger scales, one may approximate the fluid as continuously distributed.

This allows physical quantities to be represented as differentiable fields.

The continuum approximation is valid when the characteristic system scale is
much larger than the molecular mean free path.

%%%%%%%%%%%%%%%%%%%%%%%%%%%%%%%%%%%%%%%%%%%%%%%%%%%%%%%%%%%%
\subsection{Eulerian Description}
\label{subsec:fluid-eulerian-description}
%%%%%%%%%%%%%%%%%%%%%%%%%%%%%%%%%%%%%%%%%%%%%%%%%%%%%%%%%%%%

Fluid dynamics most commonly uses the Eulerian description.

One observes the fluid at fixed spatial positions.

For example,

\[
\boxed{
\mathbf{v}
=
\mathbf{v}(\mathbf{x},t)
}
\]

gives the velocity of fluid currently occupying location \(\mathbf{x}\).

This contrasts with following individual material particles.

%%%%%%%%%%%%%%%%%%%%%%%%%%%%%%%%%%%%%%%%%%%%%%%%%%%%%%%%%%%%
\subsection{Lagrangian Description}
\label{subsec:fluid-lagrangian-description}
%%%%%%%%%%%%%%%%%%%%%%%%%%%%%%%%%%%%%%%%%%%%%%%%%%%%%%%%%%%%

A Lagrangian description tracks individual fluid particles.

If a particle begins at label \(\mathbf{X}\), its trajectory is

\[
\boxed{
\mathbf{x}
=
\boldsymbol{\chi}(\mathbf{X},t).
}
\]

Its velocity is

\[
\boxed{
\frac{\partial\boldsymbol{\chi}}
{\partial t}.
}
\]

Particle tracking is useful for transport, mixing, and numerical simulation.

%%%%%%%%%%%%%%%%%%%%%%%%%%%%%%%%%%%%%%%%%%%%%%%%%%%%%%%%%%%%
\subsection{Velocity Field}
\label{subsec:fluid-velocity-field}
%%%%%%%%%%%%%%%%%%%%%%%%%%%%%%%%%%%%%%%%%%%%%%%%%%%%%%%%%%%%

The velocity field is

\[
\boxed{
\mathbf{v}
=
(u,v,w).
}
\]

In Cartesian coordinates,

\[
u
=
u(x,y,z,t),
\]

\[
v
=
v(x,y,z,t),
\]

\[
w
=
w(x,y,z,t).
\]

The field describes the local direction and speed of fluid motion.

%%%%%%%%%%%%%%%%%%%%%%%%%%%%%%%%%%%%%%%%%%%%%%%%%%%%%%%%%%%%
\subsection{Steady and Unsteady Flow}
\label{subsec:steady-unsteady-flow}
%%%%%%%%%%%%%%%%%%%%%%%%%%%%%%%%%%%%%%%%%%%%%%%%%%%%%%%%%%%%

A flow is steady when its Eulerian fields do not depend explicitly on time.

For velocity,

\[
\boxed{
\frac{\partial\mathbf{v}}{\partial t}
=
0.
}
\]

A flow can be steady while individual fluid particles still accelerate due
to spatial variations in velocity.

%%%%%%%%%%%%%%%%%%%%%%%%%%%%%%%%%%%%%%%%%%%%%%%%%%%%%%%%%%%%
\subsection{Material Derivative in a Fluid}
\label{subsec:fluid-material-derivative}
%%%%%%%%%%%%%%%%%%%%%%%%%%%%%%%%%%%%%%%%%%%%%%%%%%%%%%%%%%%%

For scalar field \(f(\mathbf{x},t)\),

\[
\boxed{
\frac{Df}{Dt}
=
\frac{\partial f}{\partial t}
+
\mathbf{v}\cdot\nabla f.
}
\]

For the velocity field itself,

\[
\boxed{
\frac{D\mathbf{v}}{Dt}
=
\frac{\partial\mathbf{v}}{\partial t}
+
(\mathbf{v}\cdot\nabla)\mathbf{v}.
}
\]

This is the acceleration of a fluid particle.

%%%%%%%%%%%%%%%%%%%%%%%%%%%%%%%%%%%%%%%%%%%%%%%%%%%%%%%%%%%%
\subsection{Local and Convective Acceleration}
\label{subsec:local-convective-acceleration}
%%%%%%%%%%%%%%%%%%%%%%%%%%%%%%%%%%%%%%%%%%%%%%%%%%%%%%%%%%%%

Fluid acceleration decomposes as

\[
\boxed{
\mathbf{a}
=
\frac{\partial\mathbf{v}}{\partial t}
+
(\mathbf{v}\cdot\nabla)\mathbf{v}.
}
\]

The first term is local acceleration.

The second is convective acceleration.

Thus a steady flow can still have nonzero acceleration if velocity changes
from place to place.

%%%%%%%%%%%%%%%%%%%%%%%%%%%%%%%%%%%%%%%%%%%%%%%%%%%%%%%%%%%%
\subsection{Worked Example: Convective Acceleration}
\label{subsec:worked-convective-acceleration}
%%%%%%%%%%%%%%%%%%%%%%%%%%%%%%%%%%%%%%%%%%%%%%%%%%%%%%%%%%%%

\begin{workedexamplebox}{Steady One-Dimensional Flow}

Suppose

\[
v(x)=ax
\]

in one dimension.

Since the flow is steady,

\[
\frac{\partial v}{\partial t}=0.
\]

The material acceleration is

\[
\frac{Dv}{Dt}
=
v\frac{dv}{dx}.
\]

Because

\[
\frac{dv}{dx}=a,
\]

we obtain

\[
\frac{Dv}{Dt}
=
(ax)a.
\]

Thus

\begin{answerbox}
\[
\boxed{
\frac{Dv}{Dt}
=
a^2x.
}
\]
\end{answerbox}

The flow accelerates despite having no explicit time dependence.

\end{workedexamplebox}

%%%%%%%%%%%%%%%%%%%%%%%%%%%%%%%%%%%%%%%%%%%%%%%%%%%%%%%%%%%%
\subsection{Streamlines}
\label{subsec:streamlines}
%%%%%%%%%%%%%%%%%%%%%%%%%%%%%%%%%%%%%%%%%%%%%%%%%%%%%%%%%%%%

A streamline is a curve tangent everywhere to the instantaneous velocity
field.

If

\[
\mathbf{x}(s)
\]

parameterizes the curve, then

\[
\boxed{
\frac{d\mathbf{x}}{ds}
\parallel
\mathbf{v}.
}
\]

In two dimensions,

\[
\boxed{
\frac{dy}{dx}
=
\frac{v}{u},
}
\]

when \(u\neq0\).

%%%%%%%%%%%%%%%%%%%%%%%%%%%%%%%%%%%%%%%%%%%%%%%%%%%%%%%%%%%%
\subsection{Pathlines}
\label{subsec:pathlines}
%%%%%%%%%%%%%%%%%%%%%%%%%%%%%%%%%%%%%%%%%%%%%%%%%%%%%%%%%%%%

A pathline is the actual trajectory of one fluid particle.

It satisfies

\[
\boxed{
\frac{d\mathbf{x}}{dt}
=
\mathbf{v}(\mathbf{x},t).
}
\]

In steady flow, streamlines and pathlines coincide.

In unsteady flow, they generally differ.

%%%%%%%%%%%%%%%%%%%%%%%%%%%%%%%%%%%%%%%%%%%%%%%%%%%%%%%%%%%%
\subsection{Streaklines}
\label{subsec:streaklines}
%%%%%%%%%%%%%%%%%%%%%%%%%%%%%%%%%%%%%%%%%%%%%%%%%%%%%%%%%%%%

A streakline is the set of particles that have previously passed through a
specified spatial point.

Smoke released continuously from a fixed source visualizes a streakline.

In steady flow,

\[
\boxed{
\text{streamlines}
=
\text{pathlines}
=
\text{streaklines}.
}
\]

%%%%%%%%%%%%%%%%%%%%%%%%%%%%%%%%%%%%%%%%%%%%%%%%%%%%%%%%%%%%
\subsection{Fluid Density}
\label{subsec:fluid-density}
%%%%%%%%%%%%%%%%%%%%%%%%%%%%%%%%%%%%%%%%%%%%%%%%%%%%%%%%%%%%

Density is mass per unit volume:

\[
\boxed{
\rho
=
\frac{dm}{dV}.
}
\]

In general,

\[
\rho
=
\rho(\mathbf{x},t).
\]

Liquids are often approximated as incompressible.

Gases frequently require variable-density models.

%%%%%%%%%%%%%%%%%%%%%%%%%%%%%%%%%%%%%%%%%%%%%%%%%%%%%%%%%%%%
\subsection{Conservation of Mass}
\label{subsec:fluid-mass-conservation}
%%%%%%%%%%%%%%%%%%%%%%%%%%%%%%%%%%%%%%%%%%%%%%%%%%%%%%%%%%%%

Mass conservation gives

\[
\boxed{
\frac{\partial\rho}{\partial t}
+
\nabla\cdot(\rho\mathbf{v})
=
0.
}
\]

This is the continuity equation.

Expanding the divergence,

\[
\boxed{
\frac{D\rho}{Dt}
+
\rho\nabla\cdot\mathbf{v}
=
0.
}
\]

%%%%%%%%%%%%%%%%%%%%%%%%%%%%%%%%%%%%%%%%%%%%%%%%%%%%%%%%%%%%
\subsection{Incompressible Flow}
\label{subsec:incompressible-flow}
%%%%%%%%%%%%%%%%%%%%%%%%%%%%%%%%%%%%%%%%%%%%%%%%%%%%%%%%%%%%

For constant density,

\[
\frac{D\rho}{Dt}=0.
\]

Therefore mass conservation becomes

\[
\boxed{
\nabla\cdot\mathbf{v}
=
0.
}
\]

This is the incompressibility condition.

It means local fluid volume is preserved.

%%%%%%%%%%%%%%%%%%%%%%%%%%%%%%%%%%%%%%%%%%%%%%%%%%%%%%%%%%%%
\subsection{Worked Example: Checking Incompressibility}
\label{subsec:worked-checking-incompressibility}
%%%%%%%%%%%%%%%%%%%%%%%%%%%%%%%%%%%%%%%%%%%%%%%%%%%%%%%%%%%%

\begin{workedexamplebox}{A Two-Dimensional Velocity Field}

Let

\[
\mathbf{v}(x,y)
=
\begin{bmatrix}
ax\\
-ay
\end{bmatrix}.
\]

Then

\[
\nabla\cdot\mathbf{v}
=
\frac{\partial(ax)}{\partial x}
+
\frac{\partial(-ay)}{\partial y}.
\]

Thus

\[
\nabla\cdot\mathbf{v}
=
a-a
=
0.
\]

\begin{answerbox}
\[
\boxed{
\nabla\cdot\mathbf{v}=0,
}
\]

so the flow is incompressible.
\end{answerbox}

\end{workedexamplebox}

%%%%%%%%%%%%%%%%%%%%%%%%%%%%%%%%%%%%%%%%%%%%%%%%%%%%%%%%%%%%
\subsection{Pressure}
\label{subsec:fluid-pressure}
%%%%%%%%%%%%%%%%%%%%%%%%%%%%%%%%%%%%%%%%%%%%%%%%%%%%%%%%%%%%

Pressure is the isotropic normal stress in a fluid.

For an inviscid fluid, the stress tensor is

\[
\boxed{
\boldsymbol{\sigma}
=
-p\mathbf{I}.
}
\]

The negative sign reflects the usual convention that positive pressure is
compressive.

%%%%%%%%%%%%%%%%%%%%%%%%%%%%%%%%%%%%%%%%%%%%%%%%%%%%%%%%%%%%
\subsection{Hydrostatics}
\label{subsec:hydrostatics}
%%%%%%%%%%%%%%%%%%%%%%%%%%%%%%%%%%%%%%%%%%%%%%%%%%%%%%%%%%%%

A fluid at rest satisfies

\[
\mathbf{v}=0.
\]

Momentum balance becomes

\[
\boxed{
0
=
-\nabla p
+
\rho\mathbf{b}.
}
\]

For gravity,

\[
\mathbf{b}
=
-g\mathbf{e}_z.
\]

Therefore,

\[
\boxed{
\frac{dp}{dz}
=
-\rho g.
}
\]

%%%%%%%%%%%%%%%%%%%%%%%%%%%%%%%%%%%%%%%%%%%%%%%%%%%%%%%%%%%%
\subsection{Hydrostatic Pressure}
\label{subsec:hydrostatic-pressure}
%%%%%%%%%%%%%%%%%%%%%%%%%%%%%%%%%%%%%%%%%%%%%%%%%%%%%%%%%%%%

If density is constant,

\[
\frac{dp}{dz}
=
-\rho g.
\]

Integrating,

\[
\boxed{
p(z)
=
p_0
+
\rho g(z_0-z).
}
\]

Thus pressure increases with depth.

%%%%%%%%%%%%%%%%%%%%%%%%%%%%%%%%%%%%%%%%%%%%%%%%%%%%%%%%%%%%
\subsection{Worked Example: Pressure With Depth}
\label{subsec:worked-pressure-depth}
%%%%%%%%%%%%%%%%%%%%%%%%%%%%%%%%%%%%%%%%%%%%%%%%%%%%%%%%%%%%

\begin{workedexamplebox}{Water Pressure}

Suppose water has density

\[
\rho=1000\text{ kg/m}^3
\]

and depth is

\[
h=5\text{ m}.
\]

The gauge pressure increase is

\[
\Delta p
=
\rho gh.
\]

Using

\[
g=9.81\text{ m/s}^2,
\]

we obtain

\[
\Delta p
=
1000(9.81)(5).
\]

Thus

\begin{answerbox}
\[
\boxed{
\Delta p
=
4.905\times10^4\text{ Pa}.
}
\]
\end{answerbox}

\end{workedexamplebox}

%%%%%%%%%%%%%%%%%%%%%%%%%%%%%%%%%%%%%%%%%%%%%%%%%%%%%%%%%%%%
\subsection{Viscosity}
\label{subsec:fluid-viscosity}
%%%%%%%%%%%%%%%%%%%%%%%%%%%%%%%%%%%%%%%%%%%%%%%%%%%%%%%%%%%%

Viscosity measures resistance to relative motion between neighboring fluid
layers.

For a simple Newtonian shear flow,

\[
\boxed{
\tau
=
\mu
\frac{du}{dy},
}
\]

where

\[
\mu
\]

is dynamic viscosity.

Large \(\mu\) corresponds to strong viscous resistance.

%%%%%%%%%%%%%%%%%%%%%%%%%%%%%%%%%%%%%%%%%%%%%%%%%%%%%%%%%%%%
\subsection{Newtonian Fluids}
\label{subsec:newtonian-fluids}
%%%%%%%%%%%%%%%%%%%%%%%%%%%%%%%%%%%%%%%%%%%%%%%%%%%%%%%%%%%%

A Newtonian fluid has viscous stress proportional to the rate of deformation.

For an incompressible Newtonian fluid,

\[
\boxed{
\boldsymbol{\sigma}
=
-p\mathbf{I}
+
2\mu\mathbf{D},
}
\]

where

\[
\boxed{
\mathbf{D}
=
\frac12
\left(
\nabla\mathbf{v}
+
(\nabla\mathbf{v})^T
\right).
}
\]

%%%%%%%%%%%%%%%%%%%%%%%%%%%%%%%%%%%%%%%%%%%%%%%%%%%%%%%%%%%%
\subsection{General Newtonian Stress Law}
\label{subsec:general-newtonian-stress}
%%%%%%%%%%%%%%%%%%%%%%%%%%%%%%%%%%%%%%%%%%%%%%%%%%%%%%%%%%%%

For a compressible Newtonian fluid, one often writes

\[
\boxed{
\boldsymbol{\sigma}
=
-p\mathbf{I}
+
2\mu\mathbf{D}
+
\lambda_v
(\nabla\cdot\mathbf{v})\mathbf{I},
}
\]

where \(\lambda_v\) is a second viscosity coefficient.

For incompressible flow,

\[
\nabla\cdot\mathbf{v}=0,
\]

so the final term vanishes.

%%%%%%%%%%%%%%%%%%%%%%%%%%%%%%%%%%%%%%%%%%%%%%%%%%%%%%%%%%%%
\subsection{Kinematic Viscosity}
\label{subsec:kinematic-viscosity}
%%%%%%%%%%%%%%%%%%%%%%%%%%%%%%%%%%%%%%%%%%%%%%%%%%%%%%%%%%%%

Kinematic viscosity is defined by

\[
\boxed{
\nu
=
\frac{\mu}{\rho}.
}
\]

Its dimensions are

\[
\boxed{
[\nu]
=
L^2T^{-1}.
}
\]

This has the same dimensions as a diffusion coefficient.

Indeed, viscosity diffuses momentum through a fluid.

%%%%%%%%%%%%%%%%%%%%%%%%%%%%%%%%%%%%%%%%%%%%%%%%%%%%%%%%%%%%
\subsection{Balance of Linear Momentum}
\label{subsec:fluid-linear-momentum}
%%%%%%%%%%%%%%%%%%%%%%%%%%%%%%%%%%%%%%%%%%%%%%%%%%%%%%%%%%%%

The continuum momentum equation is

\[
\boxed{
\rho
\frac{D\mathbf{v}}{Dt}
=
\nabla\cdot\boldsymbol{\sigma}
+
\rho\mathbf{b}.
}
\]

This equation becomes the Euler or Navier--Stokes equations depending on the
chosen stress law.

%%%%%%%%%%%%%%%%%%%%%%%%%%%%%%%%%%%%%%%%%%%%%%%%%%%%%%%%%%%%
\subsection{Euler Equations}
\label{subsec:euler-fluid-equations}
%%%%%%%%%%%%%%%%%%%%%%%%%%%%%%%%%%%%%%%%%%%%%%%%%%%%%%%%%%%%

For an inviscid fluid,

\[
\boldsymbol{\sigma}
=
-p\mathbf{I}.
\]

Therefore,

\[
\nabla\cdot\boldsymbol{\sigma}
=
-\nabla p.
\]

Momentum balance becomes

\[
\boxed{
\rho
\left[
\frac{\partial\mathbf{v}}{\partial t}
+
(\mathbf{v}\cdot\nabla)\mathbf{v}
\right]
=
-\nabla p
+
\rho\mathbf{b}.
}
\]

These are the Euler equations.

%%%%%%%%%%%%%%%%%%%%%%%%%%%%%%%%%%%%%%%%%%%%%%%%%%%%%%%%%%%%
\subsection{Navier--Stokes Equations}
\label{subsec:navier-stokes-equations}
%%%%%%%%%%%%%%%%%%%%%%%%%%%%%%%%%%%%%%%%%%%%%%%%%%%%%%%%%%%%

For a Newtonian fluid with constant viscosity and incompressible flow,

\[
\boxed{
\rho
\left[
\frac{\partial\mathbf{v}}{\partial t}
+
(\mathbf{v}\cdot\nabla)\mathbf{v}
\right]
=
-\nabla p
+
\mu\nabla^2\mathbf{v}
+
\rho\mathbf{b}.
}
\]

Together with

\[
\boxed{
\nabla\cdot\mathbf{v}=0,
}
\]

these are the incompressible Navier--Stokes equations.

%%%%%%%%%%%%%%%%%%%%%%%%%%%%%%%%%%%%%%%%%%%%%%%%%%%%%%%%%%%%
\subsection{Meaning of the Navier--Stokes Terms}
\label{subsec:navier-stokes-term-meaning}
%%%%%%%%%%%%%%%%%%%%%%%%%%%%%%%%%%%%%%%%%%%%%%%%%%%%%%%%%%%%

The term

\[
\rho\frac{\partial\mathbf{v}}{\partial t}
\]

represents local inertia.

The term

\[
\boxed{
\rho(\mathbf{v}\cdot\nabla)\mathbf{v}
}
\]

represents convective inertia.

The term

\[
\boxed{
-\nabla p
}
\]

is pressure forcing.

The term

\[
\boxed{
\mu\nabla^2\mathbf{v}
}
\]

represents viscous momentum diffusion.

%%%%%%%%%%%%%%%%%%%%%%%%%%%%%%%%%%%%%%%%%%%%%%%%%%%%%%%%%%%%
\subsection{Nonlinearity of Fluid Motion}
\label{subsec:fluid-nonlinearity}
%%%%%%%%%%%%%%%%%%%%%%%%%%%%%%%%%%%%%%%%%%%%%%%%%%%%%%%%%%%%

The term

\[
\boxed{
(\mathbf{v}\cdot\nabla)\mathbf{v}
}
\]

is nonlinear in velocity.

This nonlinearity allows fluid systems to exhibit:

\[
\boxed{
\text{instability},
}
\]

\[
\boxed{
\text{vortex interaction},
}
\]

\[
\boxed{
\text{chaotic dynamics},
}
\]

and

\[
\boxed{
\text{turbulence}.
}
\]

%%%%%%%%%%%%%%%%%%%%%%%%%%%%%%%%%%%%%%%%%%%%%%%%%%%%%%%%%%%%
\subsection{No-Slip Boundary Condition}
\label{subsec:no-slip-condition}
%%%%%%%%%%%%%%%%%%%%%%%%%%%%%%%%%%%%%%%%%%%%%%%%%%%%%%%%%%%%

For a viscous fluid touching a stationary solid wall, the standard no-slip
condition is

\[
\boxed{
\mathbf{v}=0.
}
\]

For a moving wall with velocity \(\mathbf{V}_w\),

\[
\boxed{
\mathbf{v}
=
\mathbf{V}_w.
}
\]

The fluid velocity matches the wall velocity.

%%%%%%%%%%%%%%%%%%%%%%%%%%%%%%%%%%%%%%%%%%%%%%%%%%%%%%%%%%%%
\subsection{No-Penetration Condition}
\label{subsec:no-penetration-condition}
%%%%%%%%%%%%%%%%%%%%%%%%%%%%%%%%%%%%%%%%%%%%%%%%%%%%%%%%%%%%

For an impermeable wall,

\[
\boxed{
\mathbf{v}\cdot\mathbf{n}
=
\mathbf{V}_w\cdot\mathbf{n}.
}
\]

For a stationary wall,

\[
\boxed{
\mathbf{v}\cdot\mathbf{n}=0.
}
\]

In inviscid flow, tangential slip may still be allowed.

%%%%%%%%%%%%%%%%%%%%%%%%%%%%%%%%%%%%%%%%%%%%%%%%%%%%%%%%%%%%
\subsection{Vorticity}
\label{subsec:fluid-vorticity}
%%%%%%%%%%%%%%%%%%%%%%%%%%%%%%%%%%%%%%%%%%%%%%%%%%%%%%%%%%%%

Vorticity is

\[
\boxed{
\boldsymbol{\omega}
=
\nabla\times\mathbf{v}.
}
\]

It measures local rotational motion of the fluid.

A flow is irrotational when

\[
\boxed{
\boldsymbol{\omega}=0.
}
\]

%%%%%%%%%%%%%%%%%%%%%%%%%%%%%%%%%%%%%%%%%%%%%%%%%%%%%%%%%%%%
\subsection{Worked Example: Vorticity}
\label{subsec:worked-vorticity}
%%%%%%%%%%%%%%%%%%%%%%%%%%%%%%%%%%%%%%%%%%%%%%%%%%%%%%%%%%%%

\begin{workedexamplebox}{Solid-Body Rotation}

Consider

\[
\mathbf{v}(x,y)
=
\begin{bmatrix}
-\Omega y\\
\Omega x\\
0
\end{bmatrix}.
\]

Then

\[
\nabla\times\mathbf{v}
=
\begin{bmatrix}
0\\
0\\
2\Omega
\end{bmatrix}.
\]

Thus

\begin{answerbox}
\[
\boxed{
\boldsymbol{\omega}
=
2\Omega\mathbf{e}_z.
}
\]
\end{answerbox}

Rigid rotation therefore has uniform nonzero vorticity.

\end{workedexamplebox}

%%%%%%%%%%%%%%%%%%%%%%%%%%%%%%%%%%%%%%%%%%%%%%%%%%%%%%%%%%%%
\subsection{Circulation}
\label{subsec:fluid-circulation}
%%%%%%%%%%%%%%%%%%%%%%%%%%%%%%%%%%%%%%%%%%%%%%%%%%%%%%%%%%%%

Circulation around a closed curve \(C\) is

\[
\boxed{
\Gamma
=
\oint_C
\mathbf{v}\cdot d\mathbf{r}.
}
\]

By Stokes' theorem,

\[
\boxed{
\Gamma
=
\int_S
\boldsymbol{\omega}\cdot\mathbf{n}\,dS.
}
\]

Thus circulation measures total vorticity flux through a surface bounded by
the curve.

%%%%%%%%%%%%%%%%%%%%%%%%%%%%%%%%%%%%%%%%%%%%%%%%%%%%%%%%%%%%
\subsection{Kelvin's Circulation Theorem}
\label{subsec:kelvin-circulation}
%%%%%%%%%%%%%%%%%%%%%%%%%%%%%%%%%%%%%%%%%%%%%%%%%%%%%%%%%%%%

Under suitable inviscid and conservative-force assumptions, circulation
around a closed material loop is conserved:

\[
\boxed{
\frac{D\Gamma}{Dt}
=
0.
}
\]

This theorem connects vorticity evolution with fluid-particle motion.

%%%%%%%%%%%%%%%%%%%%%%%%%%%%%%%%%%%%%%%%%%%%%%%%%%%%%%%%%%%%
\subsection{Velocity Potential}
\label{subsec:velocity-potential}
%%%%%%%%%%%%%%%%%%%%%%%%%%%%%%%%%%%%%%%%%%%%%%%%%%%%%%%%%%%%

For irrotational flow in a simply connected region,

\[
\boxed{
\nabla\times\mathbf{v}=0
}
\]

implies the existence of a velocity potential \(\phi\) such that

\[
\boxed{
\mathbf{v}
=
\nabla\phi.
}
\]

For incompressible flow,

\[
\nabla\cdot\mathbf{v}=0,
\]

so

\[
\boxed{
\nabla^2\phi=0.
}
\]

Thus potential flow reduces to Laplace's equation.

%%%%%%%%%%%%%%%%%%%%%%%%%%%%%%%%%%%%%%%%%%%%%%%%%%%%%%%%%%%%
\subsection{Stream Function}
\label{subsec:stream-function}
%%%%%%%%%%%%%%%%%%%%%%%%%%%%%%%%%%%%%%%%%%%%%%%%%%%%%%%%%%%%

For two-dimensional incompressible flow, define a stream function \(\psi\)
by

\[
\boxed{
u
=
\frac{\partial\psi}{\partial y},
\qquad
v
=
-\frac{\partial\psi}{\partial x}.
}
\]

Then

\[
\frac{\partial u}{\partial x}
+
\frac{\partial v}{\partial y}
=
0
\]

automatically.

Curves

\[
\boxed{
\psi=\text{constant}
}
\]

are streamlines.

%%%%%%%%%%%%%%%%%%%%%%%%%%%%%%%%%%%%%%%%%%%%%%%%%%%%%%%%%%%%
\subsection{Bernoulli's Equation}
\label{subsec:bernoulli-equation}
%%%%%%%%%%%%%%%%%%%%%%%%%%%%%%%%%%%%%%%%%%%%%%%%%%%%%%%%%%%%

For steady, incompressible, inviscid flow with conservative body forces,
Bernoulli's equation along a streamline is

\[
\boxed{
p
+
\frac12\rho v^2
+
\rho gz
=
\text{constant}.
}
\]

Each term has dimensions of pressure or energy per unit volume.

%%%%%%%%%%%%%%%%%%%%%%%%%%%%%%%%%%%%%%%%%%%%%%%%%%%%%%%%%%%%
\subsection{Bernoulli Head Form}
\label{subsec:bernoulli-head}
%%%%%%%%%%%%%%%%%%%%%%%%%%%%%%%%%%%%%%%%%%%%%%%%%%%%%%%%%%%%

Dividing by \(\rho g\),

\[
\boxed{
\frac{p}{\rho g}
+
\frac{v^2}{2g}
+
z
=
\text{constant}.
}
\]

The three terms are called:

\[
\boxed{
\text{pressure head},
}
\]

\[
\boxed{
\text{velocity head},
}
\]

and

\[
\boxed{
\text{elevation head}.
}
\]

%%%%%%%%%%%%%%%%%%%%%%%%%%%%%%%%%%%%%%%%%%%%%%%%%%%%%%%%%%%%
\subsection{Worked Example: Bernoulli Speed}
\label{subsec:worked-bernoulli-speed}
%%%%%%%%%%%%%%%%%%%%%%%%%%%%%%%%%%%%%%%%%%%%%%%%%%%%%%%%%%%%

\begin{workedexamplebox}{Horizontal Pressure-Speed Relation}

Consider steady inviscid flow at equal elevation.

Bernoulli gives

\[
p_1
+
\frac12\rho v_1^2
=
p_2
+
\frac12\rho v_2^2.
\]

Suppose

\[
v_1=2,
\qquad
v_2=4,
\qquad
\rho=1000.
\]

Then

\[
p_1-p_2
=
\frac12\rho
(v_2^2-v_1^2).
\]

Thus

\[
p_1-p_2
=
\frac12(1000)(16-4).
\]

Therefore,

\begin{answerbox}
\[
\boxed{
p_1-p_2
=
6000\text{ Pa}.
}
\]
\end{answerbox}

The faster region has lower pressure under these assumptions.

\end{workedexamplebox}

%%%%%%%%%%%%%%%%%%%%%%%%%%%%%%%%%%%%%%%%%%%%%%%%%%%%%%%%%%%%
\subsection{Stagnation Pressure}
\label{subsec:stagnation-pressure}
%%%%%%%%%%%%%%%%%%%%%%%%%%%%%%%%%%%%%%%%%%%%%%%%%%%%%%%%%%%%

For horizontal incompressible inviscid flow, bringing fluid reversibly to
rest gives stagnation pressure

\[
\boxed{
p_0
=
p
+
\frac12\rho v^2.
}
\]

The term

\[
\frac12\rho v^2
\]

is dynamic pressure.

%%%%%%%%%%%%%%%%%%%%%%%%%%%%%%%%%%%%%%%%%%%%%%%%%%%%%%%%%%%%
\subsection{Volume Flow Rate}
\label{subsec:volume-flow-rate}
%%%%%%%%%%%%%%%%%%%%%%%%%%%%%%%%%%%%%%%%%%%%%%%%%%%%%%%%%%%%

The volume flow rate through a surface \(A\) is

\[
\boxed{
Q
=
\int_A
\mathbf{v}\cdot\mathbf{n}\,dA.
}
\]

For uniform velocity normal to the surface,

\[
\boxed{
Q
=
vA.
}
\]

Its dimensions are

\[
\boxed{
L^3T^{-1}.
}
\]

%%%%%%%%%%%%%%%%%%%%%%%%%%%%%%%%%%%%%%%%%%%%%%%%%%%%%%%%%%%%
\subsection{Mass Flow Rate}
\label{subsec:mass-flow-rate}
%%%%%%%%%%%%%%%%%%%%%%%%%%%%%%%%%%%%%%%%%%%%%%%%%%%%%%%%%%%%

The mass flow rate is

\[
\boxed{
\dot m
=
\int_A
\rho
\mathbf{v}\cdot\mathbf{n}
\,dA.
}
\]

For constant density and uniform normal velocity,

\[
\boxed{
\dot m
=
\rho vA.
}
\]

%%%%%%%%%%%%%%%%%%%%%%%%%%%%%%%%%%%%%%%%%%%%%%%%%%%%%%%%%%%%
\subsection{One-Dimensional Continuity Relation}
\label{subsec:one-dimensional-continuity}
%%%%%%%%%%%%%%%%%%%%%%%%%%%%%%%%%%%%%%%%%%%%%%%%%%%%%%%%%%%%

For steady incompressible flow through a tube,

\[
\boxed{
A_1v_1
=
A_2v_2.
}
\]

Thus if cross-sectional area decreases, average velocity increases.

%%%%%%%%%%%%%%%%%%%%%%%%%%%%%%%%%%%%%%%%%%%%%%%%%%%%%%%%%%%%
\subsection{Worked Example: Area Change}
\label{subsec:worked-area-change}
%%%%%%%%%%%%%%%%%%%%%%%%%%%%%%%%%%%%%%%%%%%%%%%%%%%%%%%%%%%%

\begin{workedexamplebox}{Velocity Through a Narrower Pipe}

Suppose

\[
A_1=4A_2
\]

and

\[
v_1=2\text{ m/s}.
\]

Continuity gives

\[
A_1v_1
=
A_2v_2.
\]

Therefore,

\[
4A_2(2)
=
A_2v_2.
\]

Thus

\begin{answerbox}
\[
\boxed{
v_2=8\text{ m/s}.
}
\]
\end{answerbox}

\end{workedexamplebox}

%%%%%%%%%%%%%%%%%%%%%%%%%%%%%%%%%%%%%%%%%%%%%%%%%%%%%%%%%%%%
\subsection{Internal Flow}
\label{subsec:internal-flow}
%%%%%%%%%%%%%%%%%%%%%%%%%%%%%%%%%%%%%%%%%%%%%%%%%%%%%%%%%%%%

Internal flow occurs when fluid is bounded by solid surfaces, as in:

\[
\boxed{
\text{pipes},
}
\]

\[
\boxed{
\text{ducts},
}
\]

and

\[
\boxed{
\text{channels}.
}
\]

Viscous interaction with walls strongly influences the velocity distribution.

%%%%%%%%%%%%%%%%%%%%%%%%%%%%%%%%%%%%%%%%%%%%%%%%%%%%%%%%%%%%
\subsection{Plane Couette Flow}
\label{subsec:plane-couette-flow}
%%%%%%%%%%%%%%%%%%%%%%%%%%%%%%%%%%%%%%%%%%%%%%%%%%%%%%%%%%%%

Consider fluid between two parallel plates.

Let the lower plate be stationary and upper plate move with velocity \(U\).

For steady incompressible flow with no pressure gradient,

\[
\boxed{
\mu\frac{d^2u}{dy^2}=0.
}
\]

With

\[
u(0)=0,
\qquad
u(h)=U,
\]

the solution is

\[
\boxed{
u(y)
=
\frac{U}{h}y.
}
\]

%%%%%%%%%%%%%%%%%%%%%%%%%%%%%%%%%%%%%%%%%%%%%%%%%%%%%%%%%%%%
\subsection{Couette Shear Stress}
\label{subsec:couette-shear-stress}
%%%%%%%%%%%%%%%%%%%%%%%%%%%%%%%%%%%%%%%%%%%%%%%%%%%%%%%%%%%%

For

\[
u(y)=\frac{U}{h}y,
\]

the velocity gradient is

\[
\frac{du}{dy}
=
\frac{U}{h}.
\]

Therefore,

\[
\boxed{
\tau
=
\mu
\frac{U}{h}.
}
\]

The shear stress is uniform across the fluid layer.

%%%%%%%%%%%%%%%%%%%%%%%%%%%%%%%%%%%%%%%%%%%%%%%%%%%%%%%%%%%%
\subsection{Pressure-Driven Channel Flow}
\label{subsec:pressure-driven-channel-flow}
%%%%%%%%%%%%%%%%%%%%%%%%%%%%%%%%%%%%%%%%%%%%%%%%%%%%%%%%%%%%

For steady fully developed flow between stationary parallel plates,

\[
u=u(y),
\]

and the momentum equation reduces to

\[
\boxed{
0
=
-\frac{dp}{dx}
+
\mu
\frac{d^2u}{dy^2}.
}
\]

A constant pressure gradient produces a parabolic velocity profile.

%%%%%%%%%%%%%%%%%%%%%%%%%%%%%%%%%%%%%%%%%%%%%%%%%%%%%%%%%%%%
\subsection{Poiseuille Pipe Flow}
\label{subsec:poiseuille-pipe-flow}
%%%%%%%%%%%%%%%%%%%%%%%%%%%%%%%%%%%%%%%%%%%%%%%%%%%%%%%%%%%%

For steady fully developed laminar flow through a circular pipe of radius
\(R\),

\[
\boxed{
u(r)
=
-\frac{1}{4\mu}
\frac{dp}{dx}
\left(
R^2-r^2
\right).
}
\]

For flow in the positive \(x\)-direction,

\[
\frac{dp}{dx}<0.
\]

The velocity is largest at the centerline.

%%%%%%%%%%%%%%%%%%%%%%%%%%%%%%%%%%%%%%%%%%%%%%%%%%%%%%%%%%%%
\subsection{Maximum and Average Pipe Velocity}
\label{subsec:pipe-average-velocity}
%%%%%%%%%%%%%%%%%%%%%%%%%%%%%%%%%%%%%%%%%%%%%%%%%%%%%%%%%%%%

At \(r=0\),

\[
\boxed{
u_{\max}
=
-\frac{R^2}{4\mu}
\frac{dp}{dx}.
}
\]

The cross-sectional average velocity is

\[
\boxed{
\bar u
=
\frac12u_{\max}.
}
\]

Thus laminar pipe flow has a characteristic parabolic profile.

%%%%%%%%%%%%%%%%%%%%%%%%%%%%%%%%%%%%%%%%%%%%%%%%%%%%%%%%%%%%
\subsection{Hagen--Poiseuille Flow Rate}
\label{subsec:hagen-poiseuille-law}
%%%%%%%%%%%%%%%%%%%%%%%%%%%%%%%%%%%%%%%%%%%%%%%%%%%%%%%%%%%%

Integrating the velocity profile over the pipe cross section gives

\[
\boxed{
Q
=
-\frac{\pi R^4}{8\mu}
\frac{dp}{dx}.
}
\]

For a pipe of length \(L\) with pressure drop

\[
\Delta p
=
p_{\text{in}}-p_{\text{out}},
\]

\[
\boxed{
Q
=
\frac{\pi R^4}{8\mu L}
\Delta p.
}
\]

The \(R^4\) dependence makes flow rate highly sensitive to pipe radius.

%%%%%%%%%%%%%%%%%%%%%%%%%%%%%%%%%%%%%%%%%%%%%%%%%%%%%%%%%%%%
\subsection{Worked Example: Radius Sensitivity}
\label{subsec:worked-radius-sensitivity}
%%%%%%%%%%%%%%%%%%%%%%%%%%%%%%%%%%%%%%%%%%%%%%%%%%%%%%%%%%%%

\begin{workedexamplebox}{Doubling Pipe Radius}

For Poiseuille flow,

\[
Q
\propto
R^4.
\]

If the radius doubles,

\[
R_2=2R_1.
\]

Then

\[
\frac{Q_2}{Q_1}
=
\left(
\frac{R_2}{R_1}
\right)^4
=
2^4.
\]

Thus

\begin{answerbox}
\[
\boxed{
Q_2=16Q_1.
}
\]
\end{answerbox}

\end{workedexamplebox}

%%%%%%%%%%%%%%%%%%%%%%%%%%%%%%%%%%%%%%%%%%%%%%%%%%%%%%%%%%%%
\subsection{Reynolds Number}
\label{subsec:reynolds-number}
%%%%%%%%%%%%%%%%%%%%%%%%%%%%%%%%%%%%%%%%%%%%%%%%%%%%%%%%%%%%

The Reynolds number is

\[
\boxed{
\operatorname{Re}
=
\frac{\rho UL}{\mu}
=
\frac{UL}{\nu}.
}
\]

It compares inertial and viscous effects.

Large Reynolds number suggests stronger inertial effects.

Small Reynolds number indicates dominant viscous effects.

%%%%%%%%%%%%%%%%%%%%%%%%%%%%%%%%%%%%%%%%%%%%%%%%%%%%%%%%%%%%
\subsection{Deriving Reynolds Number by Scaling}
\label{subsec:reynolds-scaling}
%%%%%%%%%%%%%%%%%%%%%%%%%%%%%%%%%%%%%%%%%%%%%%%%%%%%%%%%%%%%

The inertial term scales like

\[
\boxed{
\rho
\frac{U^2}{L}.
}
\]

The viscous term scales like

\[
\boxed{
\mu
\frac{U}{L^2}.
}
\]

Their ratio is

\[
\frac{
\rho U^2/L
}{
\mu U/L^2
}
=
\frac{\rho UL}{\mu}.
\]

Therefore,

\[
\boxed{
\operatorname{Re}
=
\frac{\rho UL}{\mu}.
}
\]

%%%%%%%%%%%%%%%%%%%%%%%%%%%%%%%%%%%%%%%%%%%%%%%%%%%%%%%%%%%%
\subsection{Nondimensional Navier--Stokes}
\label{subsec:nondimensional-navier-stokes}
%%%%%%%%%%%%%%%%%%%%%%%%%%%%%%%%%%%%%%%%%%%%%%%%%%%%%%%%%%%%

Let

\[
\mathbf{x}
=
L\mathbf{x}^*,
\qquad
\mathbf{v}
=
U\mathbf{v}^*,
\qquad
t
=
\frac{L}{U}t^*.
\]

Using pressure scale

\[
\rho U^2,
\]

the dimensionless incompressible Navier--Stokes equation becomes

\[
\boxed{
\frac{\partial\mathbf{v}^*}{\partial t^*}
+
(\mathbf{v}^*\cdot\nabla^*)\mathbf{v}^*
=
-\nabla^*p^*
+
\frac{1}{\operatorname{Re}}
\nabla^{*2}\mathbf{v}^*.
}
\]

%%%%%%%%%%%%%%%%%%%%%%%%%%%%%%%%%%%%%%%%%%%%%%%%%%%%%%%%%%%%
\subsection{Low Reynolds Number Flow}
\label{subsec:low-reynolds-flow}
%%%%%%%%%%%%%%%%%%%%%%%%%%%%%%%%%%%%%%%%%%%%%%%%%%%%%%%%%%%%

When

\[
\operatorname{Re}\ll1,
\]

viscous effects dominate inertia.

The inertial terms may become negligible, giving the Stokes equations

\[
\boxed{
0
=
-\nabla p
+
\mu\nabla^2\mathbf{v}
+
\rho\mathbf{b},
}
\]

with

\[
\boxed{
\nabla\cdot\mathbf{v}=0.
}
\]

%%%%%%%%%%%%%%%%%%%%%%%%%%%%%%%%%%%%%%%%%%%%%%%%%%%%%%%%%%%%
\subsection{High Reynolds Number Flow}
\label{subsec:high-reynolds-flow}
%%%%%%%%%%%%%%%%%%%%%%%%%%%%%%%%%%%%%%%%%%%%%%%%%%%%%%%%%%%%

When

\[
\operatorname{Re}\gg1,
\]

inertia dominates in much of the flow.

Viscosity may nevertheless remain crucial near solid boundaries and inside
thin shear layers.

Thus high Reynolds number does not mean viscosity can always be ignored
everywhere.

%%%%%%%%%%%%%%%%%%%%%%%%%%%%%%%%%%%%%%%%%%%%%%%%%%%%%%%%%%%%
\subsection{Laminar Flow}
\label{subsec:laminar-flow}
%%%%%%%%%%%%%%%%%%%%%%%%%%%%%%%%%%%%%%%%%%%%%%%%%%%%%%%%%%%%

Laminar flow has relatively smooth and organized fluid motion.

Neighboring layers move with limited irregular mixing.

Examples include:

\[
\boxed{
\text{slow viscous flows}
}
\]

and

\[
\boxed{
\text{fully developed Poiseuille flow}
}
\]

under suitable conditions.

%%%%%%%%%%%%%%%%%%%%%%%%%%%%%%%%%%%%%%%%%%%%%%%%%%%%%%%%%%%%
\subsection{Turbulent Flow}
\label{subsec:turbulent-flow}
%%%%%%%%%%%%%%%%%%%%%%%%%%%%%%%%%%%%%%%%%%%%%%%%%%%%%%%%%%%%

Turbulent flow contains complex fluctuations across many spatial and temporal
scales.

Typical features include:

\[
\boxed{
\text{irregular velocity fluctuations},
}
\]

\[
\boxed{
\text{enhanced mixing},
}
\]

\[
\boxed{
\text{vortical structures},
}
\]

and

\[
\boxed{
\text{broad scale interaction}.
}
\]

Turbulence remains one of the most challenging areas of applied mathematics.

%%%%%%%%%%%%%%%%%%%%%%%%%%%%%%%%%%%%%%%%%%%%%%%%%%%%%%%%%%%%
\subsection{Transition to Turbulence}
\label{subsec:transition-to-turbulence}
%%%%%%%%%%%%%%%%%%%%%%%%%%%%%%%%%%%%%%%%%%%%%%%%%%%%%%%%%%%%

Increasing Reynolds number can destabilize laminar flows.

The transition depends not only on Reynolds number but also on:

\[
\boxed{
\text{geometry},
}
\]

\[
\boxed{
\text{disturbance amplitude},
}
\]

\[
\boxed{
\text{surface roughness},
}
\]

and

\[
\boxed{
\text{boundary conditions}.
}
\]

There is no single universal transition Reynolds number for all flows.

%%%%%%%%%%%%%%%%%%%%%%%%%%%%%%%%%%%%%%%%%%%%%%%%%%%%%%%%%%%%
\subsection{Boundary Layers}
\label{subsec:boundary-layers}
%%%%%%%%%%%%%%%%%%%%%%%%%%%%%%%%%%%%%%%%%%%%%%%%%%%%%%%%%%%%

At high Reynolds number, viscous effects are often concentrated near solid
walls.

A thin region forms in which velocity changes rapidly from the wall value to
the outer-flow value.

This region is the boundary layer.

%%%%%%%%%%%%%%%%%%%%%%%%%%%%%%%%%%%%%%%%%%%%%%%%%%%%%%%%%%%%
\subsection{Boundary-Layer Scaling}
\label{subsec:boundary-layer-scaling}
%%%%%%%%%%%%%%%%%%%%%%%%%%%%%%%%%%%%%%%%%%%%%%%%%%%%%%%%%%%%

Suppose a boundary layer has streamwise scale \(L\) and thickness \(\delta\),
with

\[
\delta\ll L.
\]

Velocity gradients normal to the wall scale like

\[
\boxed{
\frac{U}{\delta},
}
\]

which can be much larger than streamwise gradients of scale

\[
\boxed{
\frac{U}{L}.
}
\]

This makes viscosity important inside the thin layer even when
\(\operatorname{Re}\) is large.

%%%%%%%%%%%%%%%%%%%%%%%%%%%%%%%%%%%%%%%%%%%%%%%%%%%%%%%%%%%%
\subsection{Flow Separation}
\label{subsec:flow-separation}
%%%%%%%%%%%%%%%%%%%%%%%%%%%%%%%%%%%%%%%%%%%%%%%%%%%%%%%%%%%%

A boundary layer can detach from a solid surface.

This phenomenon is called flow separation.

Separation is associated with:

\[
\boxed{
\text{adverse pressure gradients},
}
\]

\[
\boxed{
\text{strong deceleration near walls},
}
\]

and

\[
\boxed{
\text{large wake formation}.
}
\]

It strongly influences drag and lift.

%%%%%%%%%%%%%%%%%%%%%%%%%%%%%%%%%%%%%%%%%%%%%%%%%%%%%%%%%%%%
\subsection{External Flow}
\label{subsec:external-flow}
%%%%%%%%%%%%%%%%%%%%%%%%%%%%%%%%%%%%%%%%%%%%%%%%%%%%%%%%%%%%

External flow occurs around bodies such as:

\[
\boxed{
\text{airfoils},
}
\]

\[
\boxed{
\text{vehicles},
}
\]

\[
\boxed{
\text{buildings},
}
\]

and

\[
\boxed{
\text{spheres}.
}
\]

Important quantities include drag, lift, separation, wakes, and boundary
layers.

%%%%%%%%%%%%%%%%%%%%%%%%%%%%%%%%%%%%%%%%%%%%%%%%%%%%%%%%%%%%
\subsection{Drag}
\label{subsec:fluid-drag}
%%%%%%%%%%%%%%%%%%%%%%%%%%%%%%%%%%%%%%%%%%%%%%%%%%%%%%%%%%%%

Drag acts approximately parallel to the incoming flow direction.

A common scaling is

\[
\boxed{
F_D
=
\frac12
\rho U^2
C_DA,
}
\]

where

\[
C_D
\]

is the drag coefficient and \(A\) is a reference area.

%%%%%%%%%%%%%%%%%%%%%%%%%%%%%%%%%%%%%%%%%%%%%%%%%%%%%%%%%%%%
\subsection{Lift}
\label{subsec:fluid-lift}
%%%%%%%%%%%%%%%%%%%%%%%%%%%%%%%%%%%%%%%%%%%%%%%%%%%%%%%%%%%%

Lift acts approximately perpendicular to the incoming flow direction.

A common representation is

\[
\boxed{
F_L
=
\frac12
\rho U^2
C_LA.
}
\]

The lift coefficient \(C_L\) depends on geometry, angle of attack, Reynolds
number, Mach number, and flow state.

%%%%%%%%%%%%%%%%%%%%%%%%%%%%%%%%%%%%%%%%%%%%%%%%%%%%%%%%%%%%
\subsection{Circulation and Lift}
\label{subsec:circulation-and-lift}
%%%%%%%%%%%%%%%%%%%%%%%%%%%%%%%%%%%%%%%%%%%%%%%%%%%%%%%%%%%%

In idealized two-dimensional inviscid flow, circulation around a body can
produce lift.

The Kutta--Joukowski relation is

\[
\boxed{
L'
=
\rho U_\infty\Gamma,
}
\]

where \(L'\) is lift per unit span and \(\Gamma\) is circulation, under the
standard sign convention appropriate to the setup.

This connects global circulation with aerodynamic force.

%%%%%%%%%%%%%%%%%%%%%%%%%%%%%%%%%%%%%%%%%%%%%%%%%%%%%%%%%%%%
\subsection{Potential Flow}
\label{subsec:potential-flow}
%%%%%%%%%%%%%%%%%%%%%%%%%%%%%%%%%%%%%%%%%%%%%%%%%%%%%%%%%%%%

Potential flow assumes

\[
\boxed{
\nabla\times\mathbf{v}=0.
}
\]

For incompressible flow,

\[
\mathbf{v}
=
\nabla\phi
\]

and

\[
\boxed{
\nabla^2\phi=0.
}
\]

Potential flow is mathematically elegant and useful away from strongly
viscous regions.

%%%%%%%%%%%%%%%%%%%%%%%%%%%%%%%%%%%%%%%%%%%%%%%%%%%%%%%%%%%%
\subsection{D'Alembert Paradox}
\label{subsec:dalembert-paradox}
%%%%%%%%%%%%%%%%%%%%%%%%%%%%%%%%%%%%%%%%%%%%%%%%%%%%%%%%%%%%

Ideal steady incompressible inviscid potential flow around a symmetric body
can predict zero net drag.

Real bodies experience drag because viscosity creates boundary layers,
separation, and wakes.

This mismatch is known as D'Alembert's paradox.

It illustrates the danger of neglecting a physically small term that has
large effects near boundaries.

%%%%%%%%%%%%%%%%%%%%%%%%%%%%%%%%%%%%%%%%%%%%%%%%%%%%%%%%%%%%
\subsection{Compressible Flow}
\label{subsec:compressible-flow}
%%%%%%%%%%%%%%%%%%%%%%%%%%%%%%%%%%%%%%%%%%%%%%%%%%%%%%%%%%%%

In compressible flow,

\[
\boxed{
\rho
}
\]

can vary substantially.

Mass conservation remains

\[
\boxed{
\rho_t
+
\nabla\cdot(\rho\mathbf{v})
=
0.
}
\]

Momentum and energy equations must be coupled with an equation of state.

%%%%%%%%%%%%%%%%%%%%%%%%%%%%%%%%%%%%%%%%%%%%%%%%%%%%%%%%%%%%
\subsection{Equation of State}
\label{subsec:fluid-equation-of-state}
%%%%%%%%%%%%%%%%%%%%%%%%%%%%%%%%%%%%%%%%%%%%%%%%%%%%%%%%%%%%

For an ideal gas,

\[
\boxed{
p
=
\rho RT,
}
\]

where \(R\) is the specific gas constant.

This relates thermodynamic variables

\[
p,\rho,T.
\]

Compressible fluid dynamics therefore couples mechanics with thermodynamics.

%%%%%%%%%%%%%%%%%%%%%%%%%%%%%%%%%%%%%%%%%%%%%%%%%%%%%%%%%%%%
\subsection{Speed of Sound}
\label{subsec:speed-of-sound}
%%%%%%%%%%%%%%%%%%%%%%%%%%%%%%%%%%%%%%%%%%%%%%%%%%%%%%%%%%%%

The speed of sound measures propagation speed of small pressure disturbances.

For an ideal gas under adiabatic conditions,

\[
\boxed{
a
=
\sqrt{\gamma RT},
}
\]

where \(\gamma\) is the ratio of specific heats.

%%%%%%%%%%%%%%%%%%%%%%%%%%%%%%%%%%%%%%%%%%%%%%%%%%%%%%%%%%%%
\subsection{Mach Number}
\label{subsec:mach-number}
%%%%%%%%%%%%%%%%%%%%%%%%%%%%%%%%%%%%%%%%%%%%%%%%%%%%%%%%%%%%

The Mach number is

\[
\boxed{
\operatorname{Ma}
=
\frac{U}{a}.
}
\]

It compares flow speed to local sound speed.

Typical classifications are:

\[
\boxed{
\operatorname{Ma}<1
\quad
\text{subsonic},
}
\]

\[
\boxed{
\operatorname{Ma}\approx1
\quad
\text{transonic},
}
\]

and

\[
\boxed{
\operatorname{Ma}>1
\quad
\text{supersonic}.
}
\]

Very large Mach numbers are often called hypersonic.

%%%%%%%%%%%%%%%%%%%%%%%%%%%%%%%%%%%%%%%%%%%%%%%%%%%%%%%%%%%%
\subsection{Shock Waves}
\label{subsec:shock-waves}
%%%%%%%%%%%%%%%%%%%%%%%%%%%%%%%%%%%%%%%%%%%%%%%%%%%%%%%%%%%%

A shock wave is a thin region across which quantities such as

\[
\boxed{
p,
\quad
\rho,
\quad
T,
\quad
U
}
\]

change rapidly.

In idealized continuum models, the shock may be treated as a discontinuity.

Conservation laws determine jump conditions across the shock.

%%%%%%%%%%%%%%%%%%%%%%%%%%%%%%%%%%%%%%%%%%%%%%%%%%%%%%%%%%%%
\subsection{Rankine--Hugoniot Concept}
\label{subsec:rankine-hugoniot-concept}
%%%%%%%%%%%%%%%%%%%%%%%%%%%%%%%%%%%%%%%%%%%%%%%%%%%%%%%%%%%%

For a conservation law

\[
\boxed{
u_t+f(u)_x=0,
}
\]

a discontinuity moving with speed \(s\) satisfies

\[
\boxed{
s[u]
=
[f(u)],
}
\]

where

\[
[u]
=
u_R-u_L.
\]

This is a Rankine--Hugoniot jump condition.

It expresses conservation across a moving discontinuity.

%%%%%%%%%%%%%%%%%%%%%%%%%%%%%%%%%%%%%%%%%%%%%%%%%%%%%%%%%%%%
\subsection{Shallow-Water Equations}
\label{subsec:shallow-water-equations}
%%%%%%%%%%%%%%%%%%%%%%%%%%%%%%%%%%%%%%%%%%%%%%%%%%%%%%%%%%%%

When horizontal length scales greatly exceed fluid depth, vertical structure
may be averaged.

In one dimension, a basic shallow-water system is

\[
\boxed{
h_t+(hu)_x=0,
}
\]

\[
\boxed{
(hu)_t
+
\left(
hu^2+\frac12gh^2
\right)_x
=
0,
}
\]

for idealized flat-bottom inviscid flow.

Here \(h\) is fluid depth and \(u\) is depth-averaged velocity.

%%%%%%%%%%%%%%%%%%%%%%%%%%%%%%%%%%%%%%%%%%%%%%%%%%%%%%%%%%%%
\subsection{Froude Number}
\label{subsec:froude-number}
%%%%%%%%%%%%%%%%%%%%%%%%%%%%%%%%%%%%%%%%%%%%%%%%%%%%%%%%%%%%

The Froude number is

\[
\boxed{
\operatorname{Fr}
=
\frac{U}{\sqrt{gL}}.
}
\]

It compares inertial effects to gravity-wave effects.

In shallow water, one often uses depth \(h\):

\[
\boxed{
\operatorname{Fr}
=
\frac{U}{\sqrt{gh}}.
}
\]

%%%%%%%%%%%%%%%%%%%%%%%%%%%%%%%%%%%%%%%%%%%%%%%%%%%%%%%%%%%%
\subsection{Subcritical and Supercritical Flow}
\label{subsec:subcritical-supercritical}
%%%%%%%%%%%%%%%%%%%%%%%%%%%%%%%%%%%%%%%%%%%%%%%%%%%%%%%%%%%%

For shallow-water flow,

\[
\boxed{
\operatorname{Fr}<1
}
\]

is subcritical,

while

\[
\boxed{
\operatorname{Fr}>1
}
\]

is supercritical.

The distinction affects propagation of gravity-wave information and the
structure of hydraulic transitions.

%%%%%%%%%%%%%%%%%%%%%%%%%%%%%%%%%%%%%%%%%%%%%%%%%%%%%%%%%%%%
\subsection{Rotating Fluids}
\label{subsec:rotating-fluids}
%%%%%%%%%%%%%%%%%%%%%%%%%%%%%%%%%%%%%%%%%%%%%%%%%%%%%%%%%%%%

In a rotating reference frame, momentum equations include Coriolis
acceleration.

The Coriolis term has form

\[
\boxed{
2\boldsymbol{\Omega}\times\mathbf{v},
}
\]

where \(\boldsymbol{\Omega}\) is the angular velocity of the rotating frame.

This effect is important in large-scale atmospheric and oceanic flows.

%%%%%%%%%%%%%%%%%%%%%%%%%%%%%%%%%%%%%%%%%%%%%%%%%%%%%%%%%%%%
\subsection{Coriolis Parameter}
\label{subsec:coriolis-parameter}
%%%%%%%%%%%%%%%%%%%%%%%%%%%%%%%%%%%%%%%%%%%%%%%%%%%%%%%%%%%%

On a rotating planet, a simplified local Coriolis parameter is

\[
\boxed{
f
=
2\Omega\sin\phi,
}
\]

where \(\phi\) is latitude.

Large-scale horizontal momentum equations often contain terms proportional to

\[
\boxed{
f\mathbf{k}\times\mathbf{v}.
}
\]

%%%%%%%%%%%%%%%%%%%%%%%%%%%%%%%%%%%%%%%%%%%%%%%%%%%%%%%%%%%%
\subsection{Geostrophic Balance}
\label{subsec:geostrophic-balance}
%%%%%%%%%%%%%%%%%%%%%%%%%%%%%%%%%%%%%%%%%%%%%%%%%%%%%%%%%%%%

At large scales, pressure-gradient and Coriolis forces may approximately
balance:

\[
\boxed{
f\mathbf{k}\times\mathbf{v}
=
-\frac{1}{\rho}\nabla_h p.
}
\]

This is geostrophic balance.

It is fundamental in atmospheric and oceanic dynamics.

%%%%%%%%%%%%%%%%%%%%%%%%%%%%%%%%%%%%%%%%%%%%%%%%%%%%%%%%%%%%
\subsection{Scalar Transport}
\label{subsec:fluid-scalar-transport}
%%%%%%%%%%%%%%%%%%%%%%%%%%%%%%%%%%%%%%%%%%%%%%%%%%%%%%%%%%%%

A scalar concentration \(c(\mathbf{x},t)\) transported by a fluid may satisfy

\[
\boxed{
\frac{\partial c}{\partial t}
+
\mathbf{v}\cdot\nabla c
=
D\nabla^2c
+
R(c),
}
\]

where

\[
D
\]

is diffusivity and \(R(c)\) represents sources, sinks, or reactions.

%%%%%%%%%%%%%%%%%%%%%%%%%%%%%%%%%%%%%%%%%%%%%%%%%%%%%%%%%%%%
\subsection{Advection}
\label{subsec:fluid-advection}
%%%%%%%%%%%%%%%%%%%%%%%%%%%%%%%%%%%%%%%%%%%%%%%%%%%%%%%%%%%%

Pure advection satisfies

\[
\boxed{
c_t
+
\mathbf{v}\cdot\nabla c
=
0.
}
\]

The scalar is transported with the flow.

Along a particle trajectory,

\[
\boxed{
\frac{Dc}{Dt}=0.
}
\]

%%%%%%%%%%%%%%%%%%%%%%%%%%%%%%%%%%%%%%%%%%%%%%%%%%%%%%%%%%%%
\subsection{Diffusion}
\label{subsec:fluid-diffusion}
%%%%%%%%%%%%%%%%%%%%%%%%%%%%%%%%%%%%%%%%%%%%%%%%%%%%%%%%%%%%

Pure diffusion satisfies

\[
\boxed{
c_t
=
D\nabla^2c.
}
\]

Diffusion tends to smooth concentration gradients.

Advection and diffusion compete in many transport problems.

%%%%%%%%%%%%%%%%%%%%%%%%%%%%%%%%%%%%%%%%%%%%%%%%%%%%%%%%%%%%
\subsection{Peclet Number}
\label{subsec:peclet-number}
%%%%%%%%%%%%%%%%%%%%%%%%%%%%%%%%%%%%%%%%%%%%%%%%%%%%%%%%%%%%

The Peclet number is

\[
\boxed{
\operatorname{Pe}
=
\frac{UL}{D}.
}
\]

It compares advective transport to diffusive transport.

Large \(\operatorname{Pe}\) indicates advection dominance.

Small \(\operatorname{Pe}\) indicates diffusion dominance.

%%%%%%%%%%%%%%%%%%%%%%%%%%%%%%%%%%%%%%%%%%%%%%%%%%%%%%%%%%%%
\subsection{Thermal Convection}
\label{subsec:thermal-convection}
%%%%%%%%%%%%%%%%%%%%%%%%%%%%%%%%%%%%%%%%%%%%%%%%%%%%%%%%%%%%

Temperature differences can create density differences.

In a gravitational field, these density differences can produce buoyant
motion.

This leads to thermal convection.

The problem couples:

\[
\boxed{
\text{momentum},
}
\]

\[
\boxed{
\text{temperature transport},
}
\]

and

\[
\boxed{
\text{buoyancy}.
}
\]

%%%%%%%%%%%%%%%%%%%%%%%%%%%%%%%%%%%%%%%%%%%%%%%%%%%%%%%%%%%%
\subsection{Boussinesq Approximation}
\label{subsec:boussinesq-approximation}
%%%%%%%%%%%%%%%%%%%%%%%%%%%%%%%%%%%%%%%%%%%%%%%%%%%%%%%%%%%%

For small density variations, one may treat density as constant except in
the buoyancy term.

A common approximation is

\[
\boxed{
\rho
\approx
\rho_0
\left[
1-\alpha(T-T_0)
\right].
}
\]

This is the basis of the Boussinesq approximation in many natural-convection
problems.

%%%%%%%%%%%%%%%%%%%%%%%%%%%%%%%%%%%%%%%%%%%%%%%%%%%%%%%%%%%%
\subsection{Rayleigh Number}
\label{subsec:rayleigh-number}
%%%%%%%%%%%%%%%%%%%%%%%%%%%%%%%%%%%%%%%%%%%%%%%%%%%%%%%%%%%%

A common Rayleigh number is

\[
\boxed{
\operatorname{Ra}
=
\frac{
g\alpha\Delta T L^3
}{
\nu\kappa
},
}
\]

where

\[
\alpha
=
\text{thermal expansion coefficient},
\]

\[
\nu
=
\text{kinematic viscosity},
\]

and

\[
\kappa
=
\text{thermal diffusivity}.
\]

It measures the tendency of buoyancy to overcome viscous and thermal
diffusion.

%%%%%%%%%%%%%%%%%%%%%%%%%%%%%%%%%%%%%%%%%%%%%%%%%%%%%%%%%%%%
\subsection{Turbulence Averaging}
\label{subsec:turbulence-averaging}
%%%%%%%%%%%%%%%%%%%%%%%%%%%%%%%%%%%%%%%%%%%%%%%%%%%%%%%%%%%%

A turbulent velocity may be decomposed as

\[
\boxed{
\mathbf{v}
=
\overline{\mathbf{v}}
+
\mathbf{v}'.
}
\]

The mean satisfies

\[
\boxed{
\overline{\mathbf{v}'}
=
0.
}
\]

Substituting this decomposition into Navier--Stokes and averaging generates
additional fluctuation terms.

%%%%%%%%%%%%%%%%%%%%%%%%%%%%%%%%%%%%%%%%%%%%%%%%%%%%%%%%%%%%
\subsection{Reynolds Stresses}
\label{subsec:reynolds-stresses}
%%%%%%%%%%%%%%%%%%%%%%%%%%%%%%%%%%%%%%%%%%%%%%%%%%%%%%%%%%%%

Reynolds averaging introduces terms involving

\[
\boxed{
\overline{v_i'v_j'}.
}
\]

These behave like additional momentum-flux terms.

They are commonly represented through a Reynolds stress tensor.

The need to model these unknown correlations creates the turbulence closure
problem.

%%%%%%%%%%%%%%%%%%%%%%%%%%%%%%%%%%%%%%%%%%%%%%%%%%%%%%%%%%%%
\subsection{Turbulence Closure Problem}
\label{subsec:turbulence-closure}
%%%%%%%%%%%%%%%%%%%%%%%%%%%%%%%%%%%%%%%%%%%%%%%%%%%%%%%%%%%%

Averaging nonlinear equations produces new unknown moments.

Equations for those moments introduce still higher-order moments.

Thus the averaged system does not automatically close.

Additional models are required.

This is the turbulence closure problem.

%%%%%%%%%%%%%%%%%%%%%%%%%%%%%%%%%%%%%%%%%%%%%%%%%%%%%%%%%%%%
\subsection{Direct Numerical Simulation}
\label{subsec:direct-numerical-simulation}
%%%%%%%%%%%%%%%%%%%%%%%%%%%%%%%%%%%%%%%%%%%%%%%%%%%%%%%%%%%%

Direct numerical simulation attempts to resolve all dynamically relevant
fluid scales without a turbulence model.

This can be extremely computationally expensive because turbulent flows
contain a wide range of scales.

DNS is therefore usually limited to moderate Reynolds numbers or relatively
simple geometries.

%%%%%%%%%%%%%%%%%%%%%%%%%%%%%%%%%%%%%%%%%%%%%%%%%%%%%%%%%%%%
\subsection{Large-Eddy Simulation}
\label{subsec:large-eddy-simulation}
%%%%%%%%%%%%%%%%%%%%%%%%%%%%%%%%%%%%%%%%%%%%%%%%%%%%%%%%%%%%

Large-eddy simulation resolves large turbulent structures while modeling
smaller unresolved scales.

Conceptually,

\[
\boxed{
\text{large scales}
=
\text{resolved},
}
\]

\[
\boxed{
\text{small scales}
=
\text{modeled}.
}
\]

This reduces computational cost compared with DNS.

%%%%%%%%%%%%%%%%%%%%%%%%%%%%%%%%%%%%%%%%%%%%%%%%%%%%%%%%%%%%
\subsection{Reynolds-Averaged Models}
\label{subsec:reynolds-averaged-models}
%%%%%%%%%%%%%%%%%%%%%%%%%%%%%%%%%%%%%%%%%%%%%%%%%%%%%%%%%%%%

Reynolds-averaged Navier--Stokes methods model mean flow rather than every
instantaneous turbulent fluctuation.

They introduce closure relations for Reynolds stresses.

These models are widely used in engineering because they are much cheaper
than resolving all turbulent scales.

%%%%%%%%%%%%%%%%%%%%%%%%%%%%%%%%%%%%%%%%%%%%%%%%%%%%%%%%%%%%
\subsection{Numerical Fluid Dynamics}
\label{subsec:numerical-fluid-dynamics}
%%%%%%%%%%%%%%%%%%%%%%%%%%%%%%%%%%%%%%%%%%%%%%%%%%%%%%%%%%%%

Computational fluid dynamics numerically approximates fluid equations.

Common methods include:

\[
\boxed{
\text{finite difference methods},
}
\]

\[
\boxed{
\text{finite volume methods},
}
\]

\[
\boxed{
\text{finite element methods},
}
\]

and

\[
\boxed{
\text{spectral methods}.
}
\]

Different formulations emphasize different mathematical properties.

%%%%%%%%%%%%%%%%%%%%%%%%%%%%%%%%%%%%%%%%%%%%%%%%%%%%%%%%%%%%
\subsection{Finite Volume Viewpoint}
\label{subsec:finite-volume-viewpoint}
%%%%%%%%%%%%%%%%%%%%%%%%%%%%%%%%%%%%%%%%%%%%%%%%%%%%%%%%%%%%

Finite volume methods integrate conservation laws over control volumes.

For a conservation law,

\[
\boxed{
u_t+\nabla\cdot\mathbf{F}(u)=0,
}
\]

integration over a cell gives

\[
\boxed{
\frac{d}{dt}
\int_Vu\,dV
=
-
\int_{\partial V}
\mathbf{F}\cdot\mathbf{n}\,dA.
}
\]

Thus conservation is built directly into the numerical formulation.

%%%%%%%%%%%%%%%%%%%%%%%%%%%%%%%%%%%%%%%%%%%%%%%%%%%%%%%%%%%%
\subsection{Numerical Stability}
\label{subsec:fluid-numerical-stability}
%%%%%%%%%%%%%%%%%%%%%%%%%%%%%%%%%%%%%%%%%%%%%%%%%%%%%%%%%%%%

A numerical scheme may become unstable if the time step is too large relative
to the spatial mesh and signal speed.

A typical explicit condition has form

\[
\boxed{
\frac{U\Delta t}{\Delta x}
\leq
C,
}
\]

where \(C\) depends on the method.

This is a CFL-type restriction.

%%%%%%%%%%%%%%%%%%%%%%%%%%%%%%%%%%%%%%%%%%%%%%%%%%%%%%%%%%%%
\subsection{Numerical Diffusion}
\label{subsec:numerical-diffusion}
%%%%%%%%%%%%%%%%%%%%%%%%%%%%%%%%%%%%%%%%%%%%%%%%%%%%%%%%%%%%

Numerical schemes may artificially smooth sharp gradients.

This effect is called numerical diffusion.

It can reduce oscillations but may also smear:

\[
\boxed{
\text{shocks},
}
\]

\[
\boxed{
\text{interfaces},
}
\]

and

\[
\boxed{
\text{thin boundary layers}.
}
\]

%%%%%%%%%%%%%%%%%%%%%%%%%%%%%%%%%%%%%%%%%%%%%%%%%%%%%%%%%%%%
\subsection{Numerical Dispersion}
\label{subsec:numerical-dispersion}
%%%%%%%%%%%%%%%%%%%%%%%%%%%%%%%%%%%%%%%%%%%%%%%%%%%%%%%%%%%%

Numerical dispersion occurs when different numerical wave components travel
at incorrect phase speeds.

This can create artificial oscillations or phase errors.

Wave and compressible-flow simulations must carefully control both numerical
dispersion and diffusion.

%%%%%%%%%%%%%%%%%%%%%%%%%%%%%%%%%%%%%%%%%%%%%%%%%%%%%%%%%%%%
\subsection{Verification and Validation in CFD}
\label{subsec:cfd-verification-validation}
%%%%%%%%%%%%%%%%%%%%%%%%%%%%%%%%%%%%%%%%%%%%%%%%%%%%%%%%%%%%

Verification asks:

\[
\boxed{
\text{Are the fluid equations being solved correctly?}
}
\]

Validation asks:

\[
\boxed{
\text{Do the equations adequately describe the physical flow?}
}
\]

Grid refinement tests numerical convergence.

Experimental comparison tests physical adequacy.

%%%%%%%%%%%%%%%%%%%%%%%%%%%%%%%%%%%%%%%%%%%%%%%%%%%%%%%%%%%%
\subsection{Grid Refinement}
\label{subsec:fluid-grid-refinement}
%%%%%%%%%%%%%%%%%%%%%%%%%%%%%%%%%%%%%%%%%%%%%%%%%%%%%%%%%%%%

Suppose simulations use mesh sizes

\[
h,
\qquad
\frac{h}{2},
\qquad
\frac{h}{4}.
\]

If computed quantities approach a limiting value as the grid is refined,
this provides evidence of numerical convergence.

Grid independence should be checked before interpreting small-scale flow
features.

%%%%%%%%%%%%%%%%%%%%%%%%%%%%%%%%%%%%%%%%%%%%%%%%%%%%%%%%%%%%
\subsection{Fluid Dynamics Workflow}
\label{subsec:fluid-dynamics-workflow}
%%%%%%%%%%%%%%%%%%%%%%%%%%%%%%%%%%%%%%%%%%%%%%%%%%%%%%%%%%%%

A practical fluid-modeling workflow is:

\begin{enumerate}
    \item define the fluid domain;
    \item identify characteristic length, velocity, and time scales;
    \item determine whether density variation matters;
    \item determine whether viscosity matters;
    \item determine whether temperature or concentration must be modeled;
    \item choose Eulerian or Lagrangian variables;
    \item impose mass conservation;
    \item impose momentum conservation;
    \item select a constitutive stress law;
    \item include gravity or other body forces;
    \item define an equation of state if compressibility matters;
    \item determine initial conditions;
    \item impose wall, inlet, outlet, or free-surface conditions;
    \item nondimensionalize the equations;
    \item compute Reynolds, Mach, Froude, Peclet, or other relevant numbers;
    \item identify dominant physical effects;
    \item determine whether analytical simplifications are possible;
    \item check for steady, irrotational, fully developed, or symmetric flow;
    \item determine whether turbulence modeling is necessary;
    \item choose an appropriate numerical method;
    \item perform mesh and time-step refinement;
    \item verify conservation numerically;
    \item compare with experiments or observations;
    \item quantify uncertainty;
    \item communicate model assumptions and limitations.
\end{enumerate}

%%%%%%%%%%%%%%%%%%%%%%%%%%%%%%%%%%%%%%%%%%%%%%%%%%%%%%%%%%%%
\subsection{Common Errors}
\label{subsec:fluid-dynamics-common-errors}
%%%%%%%%%%%%%%%%%%%%%%%%%%%%%%%%%%%%%%%%%%%%%%%%%%%%%%%%%%%%

\subsubsection{Confusing Steady Flow With Zero Acceleration}

Steady flow means

\[
\frac{\partial\mathbf{v}}{\partial t}=0.
\]

Acceleration may still arise from

\[
(\mathbf{v}\cdot\nabla)\mathbf{v}.
\]

%%%%%%%%%%%%%%%%%%%%%%%%%%%%%%%%%%%%%%%%%%%%%%%%%%%%%%%%%%%%
\subsubsection{Confusing Incompressible Flow With Zero Velocity Divergence Without Context}

For the standard incompressible continuum model,

\[
\nabla\cdot\mathbf{v}=0.
\]

More generally, mass conservation should be derived from the density field and
material-volume behavior rather than assumed blindly.

%%%%%%%%%%%%%%%%%%%%%%%%%%%%%%%%%%%%%%%%%%%%%%%%%%%%%%%%%%%%
\subsubsection{Assuming Constant Density for Every Fluid}

Liquids are often approximately constant-density.

Gases can undergo strong density variation, especially at significant Mach
numbers or temperature changes.

%%%%%%%%%%%%%%%%%%%%%%%%%%%%%%%%%%%%%%%%%%%%%%%%%%%%%%%%%%%%
\subsubsection{Confusing Pressure and Viscous Stress}

Pressure is the isotropic part of stress.

Viscous stresses arise from deformation rates.

A Newtonian fluid combines both.

%%%%%%%%%%%%%%%%%%%%%%%%%%%%%%%%%%%%%%%%%%%%%%%%%%%%%%%%%%%%
\subsubsection{Forgetting the Convective Term}

The nonlinear term

\[
(\mathbf{v}\cdot\nabla)\mathbf{v}
\]

is essential in most flows with spatial velocity variation.

%%%%%%%%%%%%%%%%%%%%%%%%%%%%%%%%%%%%%%%%%%%%%%%%%%%%%%%%%%%%
\subsubsection{Applying Bernoulli's Equation Without Checking Assumptions}

The standard Bernoulli relation requires appropriate conditions such as
steady, inviscid flow and conservative body forces.

In rotational flow, the constant may differ from streamline to streamline.

%%%%%%%%%%%%%%%%%%%%%%%%%%%%%%%%%%%%%%%%%%%%%%%%%%%%%%%%%%%%
\subsubsection{Treating Bernoulli's Equation as a Universal Pressure Law}

Pressure is determined by the full governing equations and boundary
conditions.

Bernoulli is a special consequence of those equations under restricted
assumptions.

%%%%%%%%%%%%%%%%%%%%%%%%%%%%%%%%%%%%%%%%%%%%%%%%%%%%%%%%%%%%
\subsubsection{Ignoring Viscosity at Walls}

Even when viscosity is small in the bulk flow, no-slip boundary conditions
can create strong viscous effects near walls.

%%%%%%%%%%%%%%%%%%%%%%%%%%%%%%%%%%%%%%%%%%%%%%%%%%%%%%%%%%%%
\subsubsection{Assuming High Reynolds Number Means Inviscid Everywhere}

High Reynolds number often produces thin viscous boundary layers.

Viscosity may control separation, drag, and wake formation even when its bulk
effect is small.

%%%%%%%%%%%%%%%%%%%%%%%%%%%%%%%%%%%%%%%%%%%%%%%%%%%%%%%%%%%%
\subsubsection{Confusing Laminar and Steady}

A laminar flow may be unsteady.

A turbulent flow may possess a statistically steady mean.

The terms describe different properties.

%%%%%%%%%%%%%%%%%%%%%%%%%%%%%%%%%%%%%%%%%%%%%%%%%%%%%%%%%%%%
\subsubsection{Assuming a Universal Transition Reynolds Number}

Transition depends on geometry and disturbance environment.

A Reynolds number threshold from one flow should not be transferred blindly
to another.

%%%%%%%%%%%%%%%%%%%%%%%%%%%%%%%%%%%%%%%%%%%%%%%%%%%%%%%%%%%%
\subsubsection{Confusing Streamlines and Pathlines in Unsteady Flow}

Streamlines describe instantaneous velocity geometry.

Pathlines describe actual particle trajectories.

They coincide only under steady-flow conditions.

%%%%%%%%%%%%%%%%%%%%%%%%%%%%%%%%%%%%%%%%%%%%%%%%%%%%%%%%%%%%
\subsubsection{Ignoring Units of Viscosity}

Dynamic viscosity and kinematic viscosity are different quantities:

\[
\mu
\neq
\nu.
\]

They satisfy

\[
\nu=\frac{\mu}{\rho}.
\]

%%%%%%%%%%%%%%%%%%%%%%%%%%%%%%%%%%%%%%%%%%%%%%%%%%%%%%%%%%%%
\subsubsection{Confusing Vorticity With Angular Velocity}

For rigid-body rotation, vorticity is twice the local angular velocity.

Thus these quantities are related but not identical.

%%%%%%%%%%%%%%%%%%%%%%%%%%%%%%%%%%%%%%%%%%%%%%%%%%%%%%%%%%%%
\subsubsection{Assuming Irrotational Flow Has No Curved Streamlines}

Irrotational means

\[
\nabla\times\mathbf{v}=0.
\]

Streamlines can still curve.

%%%%%%%%%%%%%%%%%%%%%%%%%%%%%%%%%%%%%%%%%%%%%%%%%%%%%%%%%%%%
\subsubsection{Assuming Potential Flow Predicts Real Drag Correctly}

Ideal potential flow can miss viscous drag mechanisms entirely.

Boundary-layer and wake effects must be included for realistic drag.

%%%%%%%%%%%%%%%%%%%%%%%%%%%%%%%%%%%%%%%%%%%%%%%%%%%%%%%%%%%%
\subsubsection{Ignoring Compressibility Near the Speed of Sound}

When Mach number is not small, density changes and acoustic effects can become
important.

Incompressible approximations may fail.

%%%%%%%%%%%%%%%%%%%%%%%%%%%%%%%%%%%%%%%%%%%%%%%%%%%%%%%%%%%%
\subsubsection{Treating Shock Waves as Ordinary Smooth Solutions}

Ideal shocks are discontinuities in continuum fields.

They require conservation-law jump conditions and weak-solution concepts.

%%%%%%%%%%%%%%%%%%%%%%%%%%%%%%%%%%%%%%%%%%%%%%%%%%%%%%%%%%%%
\subsubsection{Ignoring the Energy Equation in Compressible Flow}

Mass and momentum equations alone generally do not determine a compressible
thermodynamic flow.

An energy equation and equation of state are typically required.

%%%%%%%%%%%%%%%%%%%%%%%%%%%%%%%%%%%%%%%%%%%%%%%%%%%%%%%%%%%%
\subsubsection{Confusing Froude and Reynolds Numbers}

Reynolds number compares inertia with viscosity.

Froude number compares inertia with gravity-wave effects.

They describe different physical balances.

%%%%%%%%%%%%%%%%%%%%%%%%%%%%%%%%%%%%%%%%%%%%%%%%%%%%%%%%%%%%
\subsubsection{Confusing Peclet and Reynolds Numbers}

Peclet number compares advection with scalar diffusion.

Reynolds number compares inertia with momentum diffusion.

%%%%%%%%%%%%%%%%%%%%%%%%%%%%%%%%%%%%%%%%%%%%%%%%%%%%%%%%%%%%
\subsubsection{Ignoring Turbulence Closure}

Averaging Navier--Stokes introduces unknown Reynolds stresses.

A turbulence model or higher-fidelity simulation approach is needed.

%%%%%%%%%%%%%%%%%%%%%%%%%%%%%%%%%%%%%%%%%%%%%%%%%%%%%%%%%%%%
\subsubsection{Assuming CFD Output Is Automatically Physical}

A visually smooth simulation may still violate:

\[
\boxed{
\text{conservation},
}
\]

\[
\boxed{
\text{mesh independence},
}
\]

or

\[
\boxed{
\text{model validity}.
}
\]

Verification and validation remain necessary.

%%%%%%%%%%%%%%%%%%%%%%%%%%%%%%%%%%%%%%%%%%%%%%%%%%%%%%%%%%%%
\subsubsection{Using an Excessively Large Time Step}

Explicit numerical methods may become unstable when the time step violates
transport or wave-propagation stability restrictions.

%%%%%%%%%%%%%%%%%%%%%%%%%%%%%%%%%%%%%%%%%%%%%%%%%%%%%%%%%%%%
\subsection{Applications}
\label{subsec:fluid-dynamics-applications}
%%%%%%%%%%%%%%%%%%%%%%%%%%%%%%%%%%%%%%%%%%%%%%%%%%%%%%%%%%%%

\begin{applicationbox}

\textbf{Aerodynamics.}
Fluid equations describe lift, drag, boundary layers, wakes, aircraft flows,
and high-speed compressible motion.

\medskip

\textbf{Hydrodynamics.}
Ships, channels, waves, rivers, and underwater systems are modeled through
incompressible and free-surface flows.

\medskip

\textbf{Atmospheric Science.}
Large-scale fluid models describe winds, pressure systems, convection, and
rotating atmospheric dynamics.

\medskip

\textbf{Oceanography.}
Fluid models describe currents, waves, stratification, tides, and transport
in rotating oceans.

\medskip

\textbf{Biomechanics.}
Blood flow, airflow, and transport through biological systems are modeled
using viscous-fluid equations.

\medskip

\textbf{Chemical Engineering.}
Advection, diffusion, reaction, mixing, and multiphase transport are central
to process models.

\medskip

\textbf{Environmental Modeling.}
Pollutants, nutrients, sediments, and heat can be transported through air and
water using advection-diffusion models.

\medskip

\textbf{Energy Systems.}
Turbomachinery, combustion, heat transfer, wind turbines, and pipelines
depend on fluid-dynamic modeling.

\medskip

\textbf{Computational Fluid Dynamics.}
Numerical PDE methods approximate flows too complicated for closed-form
analysis.

\end{applicationbox}

%%%%%%%%%%%%%%%%%%%%%%%%%%%%%%%%%%%%%%%%%%%%%%%%%%%%%%%%%%%%
\subsection{Exercises}
\label{subsec:fluid-dynamics-exercises}
%%%%%%%%%%%%%%%%%%%%%%%%%%%%%%%%%%%%%%%%%%%%%%%%%%%%%%%%%%%%

\begin{exercise}[1]
Define a fluid velocity field.
\end{exercise}

\begin{exercise}[1]
Distinguish Eulerian and Lagrangian fluid descriptions.
\end{exercise}

\begin{exercise}[1]
Define a streamline.
\end{exercise}

\begin{exercise}[1]
Define a pathline.
\end{exercise}

\begin{exercise}[1]
Define a streakline.
\end{exercise}

\begin{exercise}[1]
Define vorticity.
\end{exercise}

\begin{exercise}[1]
Define circulation.
\end{exercise}

\begin{exercise}[1]
Define dynamic viscosity.
\end{exercise}

\begin{exercise}[1]
Define kinematic viscosity.
\end{exercise}

\begin{exercise}[1]
State the incompressible Navier--Stokes equations.
\end{exercise}

\begin{exercise}[2]
Compute the material acceleration for

\[
\mathbf{v}(x,y)
=
(x,-y).
\]
\end{exercise}

\begin{exercise}[2]
Determine whether

\[
\mathbf{v}(x,y)
=
(x,-y)
\]

is incompressible.
\end{exercise}

\begin{exercise}[2]
Determine whether

\[
\mathbf{v}(x,y)
=
(x,y)
\]

is incompressible.
\end{exercise}

\begin{exercise}[2]
Write the continuity equation for variable density.
\end{exercise}

\begin{exercise}[2]
Derive

\[
\nabla\cdot\mathbf{v}=0
\]

from the continuity equation for constant density.
\end{exercise}

\begin{exercise}[2]
Derive hydrostatic pressure for a constant-density fluid.
\end{exercise}

\begin{exercise}[2]
Find the pressure increase at depth \(h\) in a liquid of density \(\rho\).
\end{exercise}

\begin{exercise}[2]
Write the Newtonian shear relation

\[
\tau=\mu\frac{du}{dy}.
\]
\end{exercise}

\begin{exercise}[2]
Compute kinematic viscosity given \(\mu\) and \(\rho\).
\end{exercise}

\begin{exercise}[2]
State the Euler equations.
\end{exercise}

\begin{exercise}[2]
Identify each term in the incompressible Navier--Stokes equation.
\end{exercise}

\begin{exercise}[3]
Explain why Navier--Stokes is nonlinear.
\end{exercise}

\begin{exercise}[2]
State the no-slip condition.
\end{exercise}

\begin{exercise}[2]
State the no-penetration condition.
\end{exercise}

\begin{exercise}[2]
Compute the vorticity of

\[
\mathbf{v}
=
(-y,x,0).
\]
\end{exercise}

\begin{exercise}[3]
Determine whether

\[
\mathbf{v}
=
\nabla\phi
\]

is irrotational when \(\phi\) is twice continuously differentiable.
\end{exercise}

\begin{exercise}[3]
Show that incompressible potential flow satisfies Laplace's equation.
\end{exercise}

\begin{exercise}[2]
Define a two-dimensional stream function.
\end{exercise}

\begin{exercise}[3]
Show that a stream function automatically satisfies incompressibility.
\end{exercise}

\begin{exercise}[2]
State Bernoulli's equation.
\end{exercise}

\begin{exercise}[2]
Explain the assumptions behind Bernoulli's equation.
\end{exercise}

\begin{exercise}[3]
Use Bernoulli's equation to compare pressures at two points of different
speed but equal elevation.
\end{exercise}

\begin{exercise}[2]
Define stagnation pressure.
\end{exercise}

\begin{exercise}[2]
Define volume flow rate.
\end{exercise}

\begin{exercise}[2]
Define mass flow rate.
\end{exercise}

\begin{exercise}[2]
Derive

\[
A_1v_1=A_2v_2
\]

for steady incompressible one-dimensional flow.
\end{exercise}

\begin{exercise}[2]
Derive the velocity profile for plane Couette flow.
\end{exercise}

\begin{exercise}[2]
Compute Couette shear stress.
\end{exercise}

\begin{exercise}[3]
Derive the pressure-driven flow equation between parallel plates.
\end{exercise}

\begin{exercise}[3]
Derive the Poiseuille velocity profile in a circular pipe.
\end{exercise}

\begin{exercise}[3]
Derive the Hagen--Poiseuille flow-rate relation.
\end{exercise}

\begin{exercise}[3]
Explain the significance of the \(R^4\) dependence in pipe flow.
\end{exercise}

\begin{exercise}[2]
Define Reynolds number.
\end{exercise}

\begin{exercise}[3]
Derive Reynolds number by comparing inertial and viscous scales.
\end{exercise}

\begin{exercise}[3]
Nondimensionalize incompressible Navier--Stokes and identify
\(\operatorname{Re}\).
\end{exercise}

\begin{exercise}[2]
Explain low Reynolds number flow.
\end{exercise}

\begin{exercise}[2]
Explain high Reynolds number flow.
\end{exercise}

\begin{exercise}[2]
Distinguish laminar and turbulent flow.
\end{exercise}

\begin{exercise}[3]
Explain why a universal transition Reynolds number does not exist.
\end{exercise}

\begin{exercise}[2]
Define a boundary layer.
\end{exercise}

\begin{exercise}[3]
Explain why viscosity remains important inside a high-Reynolds-number
boundary layer.
\end{exercise}

\begin{exercise}[3]
Explain flow separation conceptually.
\end{exercise}

\begin{exercise}[2]
Define drag coefficient.
\end{exercise}

\begin{exercise}[2]
Define lift coefficient.
\end{exercise}

\begin{exercise}[3]
Explain how circulation can contribute to lift.
\end{exercise}

\begin{exercise}[3]
Explain D'Alembert's paradox.
\end{exercise}

\begin{exercise}[2]
Define Mach number.
\end{exercise}

\begin{exercise}[2]
Classify flows as subsonic or supersonic from a given Mach number.
\end{exercise}

\begin{exercise}[3]
Explain why compressibility becomes important as Mach number increases.
\end{exercise}

\begin{exercise}[3]
Explain a shock wave conceptually.
\end{exercise}

\begin{exercise}[3]
Derive the Rankine--Hugoniot relation for a scalar conservation law.
\end{exercise}

\begin{exercise}[2]
Write the one-dimensional shallow-water equations.
\end{exercise}

\begin{exercise}[2]
Define Froude number.
\end{exercise}

\begin{exercise}[3]
Distinguish subcritical and supercritical shallow-water flow.
\end{exercise}

\begin{exercise}[2]
Write the Coriolis acceleration term.
\end{exercise}

\begin{exercise}[3]
Explain geostrophic balance.
\end{exercise}

\begin{exercise}[2]
Write the advection-diffusion equation.
\end{exercise}

\begin{exercise}[2]
Define Peclet number.
\end{exercise}

\begin{exercise}[3]
Interpret large and small Peclet numbers.
\end{exercise}

\begin{exercise}[3]
Explain the Boussinesq approximation.
\end{exercise}

\begin{exercise}[2]
Define Rayleigh number conceptually.
\end{exercise}

\begin{exercise}[3]
Explain the competition represented by Rayleigh number.
\end{exercise}

\begin{exercise}[2]
Write the Reynolds decomposition

\[
\mathbf{v}
=
\overline{\mathbf{v}}
+
\mathbf{v}'.
\]
\end{exercise}

\begin{exercise}[3]
Explain Reynolds stresses.
\end{exercise}

\begin{exercise}[3]
Explain the turbulence closure problem.
\end{exercise}

\begin{exercise}[2]
Distinguish DNS, LES, and Reynolds-averaged modeling conceptually.
\end{exercise}

\begin{exercise}[2]
Explain the finite volume viewpoint.
\end{exercise}

\begin{exercise}[3]
Explain why finite volume methods are naturally conservative.
\end{exercise}

\begin{exercise}[3]
Explain a CFL-type stability restriction.
\end{exercise}

\begin{exercise}[3]
Distinguish numerical diffusion and physical viscosity.
\end{exercise}

\begin{exercise}[3]
Explain numerical dispersion.
\end{exercise}

\begin{exercise}[3]
Explain the difference between CFD verification and validation.
\end{exercise}

\begin{exercise}[3]
Describe a grid-refinement study.
\end{exercise}

\begin{exercise}[4]
Derive the incompressible Navier--Stokes equations from continuum momentum
balance and Newtonian stress.
\end{exercise}

\begin{exercise}[4]
Derive the vorticity equation for incompressible flow.
\end{exercise}

\begin{exercise}[4]
Study Kelvin's circulation theorem.
\end{exercise}

\begin{exercise}[4]
Study two-dimensional potential flow using complex variables.
\end{exercise}

\begin{exercise}[4]
Construct potential flows from sources, sinks, vortices, and uniform flow.
\end{exercise}

\begin{exercise}[4]
Study flow around a cylinder using potential flow.
\end{exercise}

\begin{exercise}[4]
Derive the Kutta--Joukowski lift relation conceptually.
\end{exercise}

\begin{exercise}[4]
Study boundary-layer scaling from Navier--Stokes.
\end{exercise}

\begin{exercise}[4]
Derive the classical boundary-layer equations.
\end{exercise}

\begin{exercise}[4]
Study Stokes flow past a sphere.
\end{exercise}

\begin{exercise}[4]
Study the derivation of Poiseuille flow in cylindrical coordinates.
\end{exercise}

\begin{exercise}[4]
Derive energy conservation for a compressible fluid.
\end{exercise}

\begin{exercise}[4]
Study one-dimensional compressible Euler equations.
\end{exercise}

\begin{exercise}[4]
Derive shock jump conditions from conservation laws.
\end{exercise}

\begin{exercise}[4]
Study shallow-water characteristics.
\end{exercise}

\begin{exercise}[4]
Study hydraulic jumps.
\end{exercise}

\begin{exercise}[4]
Derive geostrophic balance from rotating momentum equations.
\end{exercise}

\begin{exercise}[4]
Study thermal convection in a fluid layer.
\end{exercise}

\begin{exercise}[4]
Study the onset of Rayleigh--Bénard convection.
\end{exercise}

\begin{exercise}[4]
Derive Reynolds-averaged Navier--Stokes equations.
\end{exercise}

\begin{exercise}[4]
Study turbulence energy cascades conceptually.
\end{exercise}

\begin{exercise}[4]
Construct a finite volume discretization of a scalar advection equation.
\end{exercise}

\begin{exercise}[4]
Compare numerical schemes for advection using stability and numerical
diffusion.
\end{exercise}

\begin{exercise}[4]
Develop a numerical incompressible-flow simulation and perform a grid
refinement study.
\end{exercise}

%%%%%%%%%%%%%%%%%%%%%%%%%%%%%%%%%%%%%%%%%%%%%%%%%%%%%%%%%%%%
\subsection{Conceptual Summary}
\label{subsec:fluid-dynamics-summary}
%%%%%%%%%%%%%%%%%%%%%%%%%%%%%%%%%%%%%%%%%%%%%%%%%%%%%%%%%%%%

Fluid dynamics specializes continuum mechanics to flowing liquids and gases.

Mass conservation gives

\[
\boxed{
\frac{\partial\rho}{\partial t}
+
\nabla\cdot(\rho\mathbf{v})
=
0.
}
\]

For incompressible flow,

\[
\boxed{
\nabla\cdot\mathbf{v}=0.
}
\]

Momentum balance with an inviscid stress law gives the Euler equations:

\[
\boxed{
\rho
\frac{D\mathbf{v}}{Dt}
=
-\nabla p
+
\rho\mathbf{b}.
}
\]

For an incompressible Newtonian fluid, the governing equations become

\[
\boxed{
\rho
\left[
\mathbf{v}_t
+
(\mathbf{v}\cdot\nabla)\mathbf{v}
\right]
=
-\nabla p
+
\mu\nabla^2\mathbf{v}
+
\rho\mathbf{b}.
}
\]

Important structures include

\[
\boxed{
\boldsymbol{\omega}
=
\nabla\times\mathbf{v}
}
\]

for vorticity,

\[
\boxed{
\Gamma
=
\oint_C
\mathbf{v}\cdot d\mathbf{r}
}
\]

for circulation, and

\[
\boxed{
p+\frac12\rho v^2+\rho gz
=
\text{constant}
}
\]

for Bernoulli's relation under its appropriate assumptions.

Nondimensional analysis introduces quantities such as

\[
\boxed{
\operatorname{Re}
=
\frac{\rho UL}{\mu},
}
\]

\[
\boxed{
\operatorname{Ma}
=
\frac{U}{a},
}
\]

\[
\boxed{
\operatorname{Fr}
=
\frac{U}{\sqrt{gL}},
}
\]

and

\[
\boxed{
\operatorname{Pe}
=
\frac{UL}{D}.
}
\]

These numbers identify the physical processes dominating a flow.

Fluid dynamics therefore combines

\[
\boxed{
\text{conservation}
+
\text{constitutive behavior}
+
\text{transport}
+
\text{scaling}
+
\text{computation}
}
\]

to describe systems ranging from microscopic viscous motion to atmospheric,
oceanic, aerodynamic, and turbulent flows.

%%%%%%%%%%%%%%%%%%%%%%%%%%%%%%%%%%%%%%%%%%%%%%%%%%%%%%%%%%%%
\subsection{Section Connections}
\label{subsec:fluid-dynamics-connections}
%%%%%%%%%%%%%%%%%%%%%%%%%%%%%%%%%%%%%%%%%%%%%%%%%%%%%%%%%%%%

\chapterconnection{
Fluid Dynamics specializes the continuum framework of
Section~\ref{sec:continuum-mechanics} to liquids and gases. Mass, momentum,
and energy conservation produce the governing field equations, while
Newtonian constitutive laws introduce viscosity and pressure into the stress
tensor. Vector calculus describes vorticity, circulation, and potential
flow; partial differential equations describe waves, diffusion, shocks, and
transport; nondimensionalization reveals Reynolds, Mach, Froude, Peclet, and
other controlling parameters; and numerical mathematics provides the basis
for computational fluid dynamics. The next section, Control Theory,
develops mathematical methods for influencing dynamical systems through
feedback, stability analysis, state-space models, controllability,
observability, and optimal control.
}

%%%%%%%%%%%%%%%%%%%%%%%%%%%%%%%%%%%%%%%%%%%%%%%%%%%%%%%%%%%%
% SECTION SOURCE TRACKING
%%%%%%%%%%%%%%%%%%%%%%%%%%%%%%%%%%%%%%%%%%%%%%%%%%%%%%%%%%%%

%------------------------------------------------------------
% 13.5 Fluid Dynamics
%------------------------------------------------------------
%
% Source basis:
%   - Standard incompressible and compressible fluid mechanics.
%   - Standard continuum balance laws and Newtonian constitutive theory.
%   - Standard potential flow, viscous flow, boundary-layer,
%     transport, geophysical, turbulent, and computational fluid dynamics.
%
% Used for:
%   - continuum fluid description;
%   - Eulerian and Lagrangian viewpoints;
%   - velocity fields;
%   - streamlines, pathlines, and streaklines;
%   - material derivatives;
%   - local and convective acceleration;
%   - density;
%   - mass conservation;
%   - incompressibility;
%   - pressure;
%   - hydrostatics;
%   - viscosity;
%   - Newtonian fluids;
%   - dynamic and kinematic viscosity;
%   - Euler equations;
%   - Navier--Stokes equations;
%   - nonlinear convection;
%   - wall boundary conditions;
%   - vorticity;
%   - circulation;
%   - Kelvin circulation theorem;
%   - velocity potentials;
%   - stream functions;
%   - Bernoulli equation;
%   - stagnation pressure;
%   - flow rates;
%   - one-dimensional continuity;
%   - Couette flow;
%   - pressure-driven channel flow;
%   - Poiseuille flow;
%   - Hagen--Poiseuille law;
%   - Reynolds number;
%   - nondimensional Navier--Stokes;
%   - low and high Reynolds number flow;
%   - laminar flow;
%   - turbulent flow;
%   - transition;
%   - boundary layers;
%   - separation;
%   - external flows;
%   - lift and drag;
%   - circulation-generated lift;
%   - potential flow;
%   - D'Alembert paradox;
%   - compressible flow;
%   - equations of state;
%   - speed of sound;
%   - Mach number;
%   - shock waves;
%   - Rankine--Hugoniot relations;
%   - shallow-water equations;
%   - Froude number;
%   - rotating fluids;
%   - Coriolis effects;
%   - geostrophic balance;
%   - scalar transport;
%   - advection and diffusion;
%   - Peclet number;
%   - thermal convection;
%   - Boussinesq approximation;
%   - Rayleigh number;
%   - Reynolds averaging;
%   - Reynolds stresses;
%   - turbulence closure;
%   - DNS;
%   - LES;
%   - Reynolds-averaged models;
%   - computational fluid dynamics;
%   - finite volume methods;
%   - CFL-type restrictions;
%   - numerical diffusion and dispersion;
%   - CFD verification and validation;
%   - grid refinement;
%   - applications and exercises.
%
% Notes:
%   - Control Theory is developed next in Section 13.6.
%   - Network Science follows in Section 13.7.
%   - Mathematical Finance follows in Section 13.8.
%   - Mathematical Economics follows in Section 13.9.
%
%%%%%%%%%%%%%%%%%%%%%%%%%%%%%%%%%%%%%%%%%%%%%%%%%%%%%%%%%%%%